\documentclass[11pt]{article}

\usepackage[a4paper,margin=2.5cm]{geometry}
\usepackage{fontspec}
\setmainfont{Latin Modern Roman}
\usepackage{booktabs}
\usepackage{longtable}
\usepackage{array}
\usepackage[table]{xcolor}
\usepackage{listings}
\usepackage{hyperref}
\usepackage{parskip}
\usepackage{amsmath}
\usepackage{microtype}

\hypersetup{
  colorlinks=true,
  linkcolor=blue!50!black,
  citecolor=blue!50!black,
  urlcolor=blue!50!black
}

\lstset{
  basicstyle=\ttfamily\footnotesize,
  breaklines=true,
  columns=fullflexible,
  keepspaces=true,
  frame=single,
  backgroundcolor=\color{gray!5},
  captionpos=b
}

\newcommand{\commit}[1]{\texttt{#1}}
\newcommand{\code}[1]{\texttt{#1}}
\newcommand{\note}[1]{\textit{#1}}
\newcommand{\gainmeasured}{\textbf{Measured}}
\newcommand{\gainbounded}{\textbf{Bounded}}
\newcommand{\gainenabling}{\textbf{Enabling}}
\newcommand{\stateactive}{\textbf{Active}}
\newcommand{\stateoptional}{\textbf{Optional}}
\newcommand{\statereverted}{\textbf{Reverted}}
\newcommand{\stateopen}{\textbf{Open}}

\newcolumntype{I}{>{\raggedright\arraybackslash}p{2.1cm}}
\newcolumntype{O}{>{\raggedright\arraybackslash}p{3.6cm}}
\newcolumntype{R}{>{\raggedright\arraybackslash}p{3.6cm}}
\newcolumntype{G}{>{\raggedright\arraybackslash}p{4.4cm}}
\newcolumntype{K}{>{\raggedright\arraybackslash}p{4.6cm}}
\newcolumntype{S}{>{\raggedright\arraybackslash}p{1.7cm}}

\title{The Optimizations Behind cMVBT-OSIC: A Unified, Verified Account}
\author{Amir Tonta}
\date{12 August 2026}

\begin{document}
\maketitle

\begin{abstract}
This document consolidates every optimization-focused write-up in this repository
--- the narrative design document (\path{docs/optimizations.tex}), the tabular
inventory (\path{docs/optimization_catalog.tex}), the compact implementation guide
(\path{docs/index_optimizations.md}), and five focused microbenchmark reports ---
into one single account of the entire \texttt{cMVBT-OSIC} implementation. It is
organized neither by commit nor by document of origin, but bottom-up by
architectural layer, exactly the order a reader would need to actually understand
the system: physical page layout, then concurrency control built on top of it,
then structural maintenance, then the transaction layer, then garbage collection,
then the cross-cutting sharding and durability stories that sit above all of it.
Every optimization is explained (what changed, why, and what it cost or risked), not
merely listed. A dedicated section (\S\ref{sec:foundational}) argues explicitly for
four optimizations that are \emph{foundational}: removing any one of them would not
just cost some throughput, it would make several other optimizations described later
in this document either meaningless or unsound. Wherever the repository's own history
contains a controlled before/after measurement, that number is reported with its
workload, thread count, and caveats attached. Wherever it was practical to do so
within this session, those numbers were independently re-measured on the machine this
document was written on (\S\ref{sec:fresh}); the two do not always agree, and both are
reported, since the disagreement itself is informative. This document also surfaces
optimization attempts that were tried and reverted, and issues that remain open and
unfixed as of this writing --- an accurate account of the whole implementation has to
include the parts that did not work as well as the parts that did.
\end{abstract}

\tableofcontents

\section{Introduction and how to read this document}
\label{sec:intro}

\texttt{cMVBT} is a concurrent, in-memory, multiversion B-tree implementing the
asymptotically optimal MVBT structure of Becker et al.~\cite{becker1996asymptotically},
extended with the multiversion concurrency control protocol of
\cite{tonta2026multiversion}. On top of the tree, the system integrates
\emph{Ordered Snapshot Instant Commit} (OSIC), the visibility/durability model from
LeanStore's scalable snapshot-isolation work~\cite{alhomssi2023scalable}: commits
become visible instantly, in commit order, while the write-ahead log is hardened to
disk lazily and in batches, off the commit's critical path.

\paragraph{Roadmap.} Sections build on one another in this order:

\begin{enumerate}
  \item \textbf{Foundational optimizations} (\S\ref{sec:foundational}): four
        decisions that everything downstream either depends on for correctness or
        depends on for the optimization to matter at all.
  \item \textbf{Physical layer} (\S\ref{sec:physical}): how one page holds several
        versions of a key, and how that page is laid out in memory.
  \item \textbf{Concurrency control} (\S\ref{sec:cc}): how threads touch pages from
        \S\ref{sec:physical} without unnecessary blocking.
  \item \textbf{Structural maintenance} (\S\ref{sec:smo}): splitting and merging
        pages as they fill up or empty out.
  \item \textbf{The self-overwrite optimization} (\S\ref{sec:selfoverwrite}): bounding
        how much physical growth one transaction's repeated writes can cause --- the
        second load-bearing optimization this document singles out, one layer up
        from \S\ref{sec:foundational}.
  \item \textbf{Garbage collection} (\S\ref{sec:gc}): reclaiming pages retired by
        \S\ref{sec:smo}.
  \item \textbf{Root indexing} (\S\ref{sec:root}) and
        \textbf{sharding} (\S\ref{sec:sharding}): finding the right root, and the
        recurring pattern of partitioning shared bookkeeping state by worker.
  \item \textbf{OSIC and the WAL} (\S\ref{sec:osic}): the commit/durability layer,
        including a full accounting of the WAL performance work.
  \item \textbf{Range scans} (\S\ref{sec:scans}) and
        \textbf{workload/benchmark-harness optimizations} (\S\ref{sec:workload}).
  \item \textbf{Case study: \code{BigTreeSize}} (\S\ref{sec:bigtree}): one
        optimization followed end to end, including a benchmark methodology bug that
        was found and fixed mid-investigation.
  \item \textbf{Rejected experiments} (\S\ref{sec:rejected}) and
        \textbf{open issues} (\S\ref{sec:open}): what was tried and did not work, and
        what is still broken.
  \item \textbf{Fresh measurements} (\S\ref{sec:fresh}), a
        \textbf{combined gains table} (\S\ref{sec:combined}), a
        \textbf{timeline} (\S\ref{sec:timeline}), and the \textbf{conclusion}
        (\S\ref{sec:conclusion}).
\end{enumerate}

\paragraph{Evidence levels.} As in the earlier catalog, every claimed gain is tagged:
\gainmeasured{} (a before/after benchmark exists, with its workload and caveats),
\gainbounded{} (a countable cost --- an allocation, a lock, a retry, a scan --- is
removably bounded, but no trustworthy end-to-end percentage was recorded), or
\gainenabling{} (primarily a correctness or livelock/deadlock fix, not a claimed
speedup). Combined benchmark results are never attributed to one component in
isolation unless the repository contains a measurement that isolates it.

\section{Foundational optimizations: what everything else stands on}
\label{sec:foundational}

Reading the rest of this document commit-by-commit gives the impression of dozens of
independent, similarly-weighted improvements. They are not independent. Four design
decisions sit underneath nearly everything else described later, in the specific
sense that removing any one of them would not merely cost throughput --- it would
make a named optimization described in a later section either \emph{unsound} (its
correctness argument stops holding) or \emph{pointless} (the cost it removes stops
existing in the first place). This section names all four before the rest of the
document goes layer by layer through the system; each subsection below points
forward to exactly which later optimizations it underwrites.

\subsection{Append-only, publish-once physical pages}

Every leaf page (\path{src/mv_page_model/leaf_page.rs}) is a fixed-capacity array
that is only ever appended to. An obsolete, aborted, or deleted entry is \emph{marked}
dead via \path{VersionInfo::delete}/\path{invalidate}; it is never physically
overwritten or removed until the next structural-maintenance pass rebuilds the page
from scratch (\S\ref{sec:physical}). Internal pages follow the same rule: a
superseded child entry is never mutated, only shadowed by a newer appended entry
(\S\ref{sec:gc}).

\textbf{What this underwrites.} Optimistic lock coupling (\S\ref{sec:cc}) lets
readers inspect a page with no latch at all specifically because nothing they might
observe was ever mutated after being published --- a reader racing a writer either
sees the old, still-fully-formed state or the new one, never a half-written slot.
Garbage collection's internal-page liveness derivation (\path{live_mask()},
\S\ref{sec:gc}) is a pure function of \emph{disjoint, append-ordered intervals} ---
it has no meaning if slots could be mutated in place. And the self-overwrite
optimization's entire safety argument (\S\ref{sec:selfoverwrite}) is the claim
``nobody else could possibly be relying on the current content of this uncommitted
tuple''; that claim presupposes a notion of ``this transaction's own, exclusively
owned tuple'' that only exists because writes append rather than clobber. Take this
foundation away and every one of those three later sections would need a completely
different, almost certainly lock-based, design.

\subsection{Optimistic lock coupling and the SmartCell latch}

Readers take an \path{Acquire} version sample of a node's \path{SmartCell}
(\path{src/mv_sync/smart_cell.rs}) and validate it after use, with no latch acquired
at all; writers upgrade to an exclusive, seqlock-style write lock only on the node
they actually mutate (\S\ref{sec:cc}). This one protocol survives out of six that
were prototyped (\S\ref{sec:cc}).

\textbf{What this underwrites.} Every sharding optimization in \S\ref{sec:sharding}
--- per-worker \path{TxContext} slots, GC's per-shard dead-block index, cache-padding
those slots against false sharing --- is a fix for contention \emph{among concurrent,
lock-free readers and writers}. If reads instead took a page latch (as in the
discarded \path{LockCoupling}/\path{LHL}/\path{HL} protocols), most of that
contention would already be serialized by the latch, and sharding the bookkeeping
around it would buy nothing. The entire fence-removal chronology
(Table~\ref{tab:fences}) is similarly parasitic on OLC existing at all: those fixes
remove ordering that OLC's own atomics already subsumed, which is only a real
saving because OLC's hot path runs on every read.

\subsection{OSIC: decoupling visibility from durability}

Every transaction draws a start and a commit timestamp from one global atomic
counter; a write becomes visible to other workers the instant its commit timestamp
lands in the committing worker's append-only \path{CommitLog} --- an in-memory,
$O(1)$ step with no wait on disk I/O (\S\ref{sec:osic}).

\textbf{What this underwrites.} The self-overwrite fast path's detection rule,
\path{insertion_stamp() == stamp}, is a comparison of two \path{TxStamp} values ---
a concept (\path{(worker\_id, ts\_start)}, constant for a whole transaction) that
\emph{is} an OSIC concept; without OSIC's stamp model there is no single equality
check that means ``did I, this very transaction, write this'' at all. More broadly,
the entire WAL performance campaign in \S\ref{sec:osic} --- sharding the write
channel, removing a discarded flush ticket, right-sizing framing buffers, an
alternative lock-free backend --- is optimizing the cost of \emph{catching up
durability asynchronously}. If commit instead required a synchronous fsync
round-trip (the naive MVCC alternative OSIC replaces), every one of those WAL
micro-optimizations would be dwarfed by that single round-trip's latency by one to
two orders of magnitude, and none of them would have been worth attempting first.

\subsection{A per-worker identity substrate}

A \path{WorkerId} is assigned once per (tree, thread) and cached thread-locally
(\path{src/mv_sync/worker.rs}); it is the index into every per-worker structure the
rest of the system builds (\S\ref{sec:sharding}).

\textbf{What this underwrites.} Sharding, as a pattern, needs something to shard
\emph{along}. GC's dead-block tracker shards by \path{worker_id \% shards.len()};
\path{TxContext}'s live-snapshot and in-flight-registration slots are one array
entry per \path{WorkerId}; the thread-local \path{SnapshotCache} is keyed the same
way. None of those three later optimizations is a response to ``things are shared
and contended'' in the abstract --- they are specifically the removal of contention
on structures that used to be indexed by nothing (a single global mutex or counter)
and are now indexed by this one identity. Without a stable per-worker identity
assigned once and reused, each of those shards would need its own separate discovery
mechanism, and the uniform ``own-shard-first, steal on miss'' idiom used throughout
\S\ref{sec:sharding} would not generalize across them the way it currently does.

\subsection{A second-order case: self-overwrite}

One further optimization is worth flagging here even though it is not, by the
standard above, foundational in the same primary sense --- it depends on the four
foundations above, rather than being depended on by everything. It is nonetheless
\emph{load-bearing for a specific later section}: without the self-overwrite fast
path (\S\ref{sec:selfoverwrite}) bounding one transaction's repeated writes to a
single key at one extra physical tuple, the structural-maintenance tuning in
\S\ref{sec:smo} --- careful capacity thresholds, tiered boundary search for splits,
emptier-sibling selection for merges --- would be fighting an unbounded adversary: a
single transaction that updates one hot key enough times inside itself can, without
this fix, leave a page with no valid split boundary at all
(\S\ref{sec:selfoverwrite} explains the exact failure shape). All of \S\ref{sec:smo}
assumes it is dealing with \emph{bounded} per-key physical chains; self-overwrite is
what keeps that assumption true.

\section{Physical layer: pages and multiversion records}
\label{sec:physical}

\subsection{Append-only leaves and the version tuple}

A leaf's record array (\path{leaf_page.rs}) is append-only, as established in
\S\ref{sec:foundational}; a packed \path{len} counter tracks active vs.\ dead slots.
Each tuple carries a \path{VersionInfo}: an \path{insertion_stamp: TxStamp} (a packed
\path{(worker_id, ts_start)}, constant for the whole transaction --- not per write,
which is exactly what makes the self-overwrite check in \S\ref{sec:selfoverwrite} a
single equality test) and an optional \path{deletion_stamp}. Lookups scan
\emph{backward} (\path{rfind}) for the newest live entry, since physical append order
places the newest write last --- this backward-scan invariant is relied on by merge
(\S\ref{sec:smo}, which must not re-sort records by \path{ts_start}) and by the
point-read miss fix discussed in \S\ref{sec:scans}.

Structural maintenance decides what a rebuild keeps via \path{record_survives_gc}: a
record survives if it is live, \emph{or} its deleting transaction is still
unresolved (may still abort, and abort needs the predecessor to undelete). Readers
are not the reason old records are copied forward during a rebuild --- they instead
route through the retired source block until they finish, at which point that block
becomes reusable (\S\ref{sec:gc}).

\subsection{Structure-of-arrays leaf layout}

A leaf's physical layout is a dense key region followed by a parallel
version/payload region, not one array of complete \path{\{key, version, payload\}}
records:

\begin{lstlisting}
key_region:   [key0, key1, key2, ...]
data_region:  [(version0, payload0), (version1, payload1), ...]
\end{lstlisting}

Point lookup walks the key region backward; version metadata and payload cache
lines are touched only \emph{after} a key matches, and a range scan can test its key
range before loading the parallel data at all. For \path{u64} keys and a one-word
payload slot, one physical entry is 32 bytes (8 bytes key, 24 bytes version+payload)
--- the same per-slot footprint as the earlier array-of-structures layout, so this is
a pure locality change, not a capacity change.

\subsection{Inline validity mask}

A normal leaf has at most 123 physical slots, so two \path{AtomicU64}s (128 bits)
cover every slot's \emph{physical} validity (cleared on abort) directly inside the
page, replacing an earlier boxed-slice handle and its pointer chase. This is
explicitly \emph{not} an MVCC visibility mask --- visibility also depends on the
reading transaction's snapshot, which one universal per-slot bit cannot represent ---
it only lets lookup skip an aborted slot before loading its version header at all.
Deliberately oversized TPC-C leaf variants (\S\ref{sec:bigtree}) fall back to a
heap bitmap through a same-sized union so the normal 4KiB leaf does not pay for it.

\subsection{Inline internal-page liveness mask}

Internal pages are append-only across structural versions: replacing a child appends
one or more new intervals while the old interval stays physically present for older
snapshots. An earlier design tracked which entries were \emph{currently} live with a
separately mutated, in-place ``obsolete mark'' atomic bit on an already-published,
concurrently-read slot --- confirmed by ThreadSanitizer as a genuine race against
range-query iteration (\S\ref{sec:gc}, \commit{fe03d90}). The current design
(\path{InternalPage::live_mask()}) instead derives liveness from the invariant that
\emph{live sibling key intervals are always pairwise disjoint}: for every appended
interval, clear the bit of every older interval it overlaps, set its own bit, and
publish with the existing Release length store. This is an $O(\mathrm{FAN\_OUT}^2)$
pure-integer pass that runs only on the SMO cold path (append), so every later SMO
consumer --- merge-candidate selection, survivor counting, split copying --- reads
\path{is_slot_live(i)} as one word load and bit test instead of rebuilding a
temporary \path{Vec<bool>} from scratch each time. Historical routing deliberately
does \emph{not} consult this mask: a child superseded in the \emph{current}
partition may still be the correct route for an older snapshot, so historical reads
keep selecting via immutable per-entry versions (\S\ref{sec:scans}).

\subsection{Exact 4\,KiB allocation, payload indirection, and the hot/cold split}

The 4\,KiB target applies to the full allocated \path{OptCell<Block<...>>}, including
the OLC \path{cell_version} word and node metadata --- not just the record array.
With \path{FAN\_OUT = NUM\_RECORDS = 123}, the cell lands at exactly 4096 bytes; 124
records would cross that boundary. This single two-record trim
(\commit{b29d347}, 125\,$\to$\,123) is \gainmeasured{}: page-aligned allocations rose
from 25\% to 100\% of the time, with a 9--12\% reduction in dTLB load misses ---
concrete evidence that ``nominally full'' and ``lands cleanly on a hardware page''
are two different numbers, and only the second one is free.

A one-word \path{PayloadSlot} stores a \path{u64} or the thin YCSB row handle
inline; larger TPC-C rows go through a reference-counted (\path{triomphe::Arc})
pointer instead, so cloning a read result or copying a survivor during a split
increments a refcount rather than copying a whole row. \path{MVBTSt} itself is
split into a hot part (root, allocator, transaction context --- everything the
traversal-critical path touches) and a \path{cold: Box<_>} part (WAL handle, table
id, construction-time bounds) purely to shrink the cache footprint of the hot path.

\subsection{Measured: leaf layout, isolated}

The split (structure-of-arrays) layout's benefit was isolated in a standalone,
dependency-free microbenchmark (\path{tools/leaf_layout_bench.rs}) comparing it
against the earlier array-of-structures layout, with no tree traversal, latching,
WAL, GC, transactions, allocation, or payload cloning in the timed region. 16,384
independently boxed leaf pages per layout (a 62.5\,MiB working set each), five
update waves per key (append-ordered, matching real leaves), point lookups walking
backward, scans checking the range before visibility.

\textbf{Original result} (AMD Ryzen 9 5900X, median of 5 runs, 750\,ms/operation):

\begin{center}
\begin{tabular}{lrrr}
\toprule
Operation & AoS & Split & Speedup \\
\midrule
Point: latest visible & 30.57\,Mlookups/s & 33.17\,Mlookups/s & 1.085$\times$ \\
Point: 1-version fallback & 18.43\,Mlookups/s & 20.15\,Mlookups/s & 1.093$\times$ \\
Point: 3-version fallback & 9.03\,Mlookups/s & 11.02\,Mlookups/s & 1.221$\times$ \\
Scan: full leaves & 322.84\,Mslots/s & 320.65\,Mslots/s & 0.993$\times$ \\
Scan: 20\% key ranges & 354.82\,Mslots/s & 389.09\,Mslots/s & 1.097$\times$ \\
\bottomrule
\end{tabular}
\end{center}

The interpretation offered at the time: point lookups improve because the key
region can reject non-matching entries without touching the version/payload cache
lines at all, and the effect grows with how many dead versions a lookup has to walk
past (8.5\%, 9.3\%, 22.1\% for 0/1/3 fallbacks); full scans are unaffected because
they touch every byte of the leaf regardless of layout.

\textbf{Independent reproduction, this session (\S\ref{sec:fresh} has the full
protocol and discussion):} on a different, older CPU (AMD Ryzen 7 2700X, Zen+
rather than Zen 3), five repeated full-binary runs of the identical benchmark were
extremely internally consistent (well under 1\% run-to-run variance) but told a
\emph{smaller and partly reversed} story: the 62.5\,MiB multi-leaf point-lookup
gains replicated in direction but at roughly a third of the claimed magnitude
(1.03$\times$--1.08$\times$, not 1.09$\times$--1.22$\times$), and two of three scan
variants came out \emph{slower} for the split layout rather than faster
(0.93$\times$--0.97$\times$). This is very likely a real hardware effect, not a
methodology error --- see \S\ref{sec:fresh} for the reasoning --- and it is reported
here precisely because a document claiming to cover the whole implementation should
not silently drop a result that failed to reproduce as cleanly as hoped.

\section{Concurrency control: from lock coupling to optimistic lock coupling}
\label{sec:cc}

\subsection{The protocols that were tried}

A pre-multiversion sister crate, \texttt{CCBPlusTree}, defines six lock-mode
variants: \path{MonoWriter} (no real concurrency control), \path{LockCoupling}
(classic hand-over-hand exclusive latching), \path{ORWC} (``Optimistic
Read-Write Coupling''), \path{OLC} (Leis-style optimistic lock coupling), and two
level-bounded hybrids, \path{LHL}/\path{HL}. Only \textbf{OLC} survives in the
current codebase's write/read fast paths (\path{traversal_write_olc},
\path{traverse_read_key}, \path{try_retire}, \path{upgrade_write_lock}).
\path{MonoWriter}/\path{LockCoupling}/\path{LHL}/\path{HL} were prototyped in early
2024 but never carried into the multiversion rewrite.

\subsection{The rise and fall of ORWC --- a cautionary tale}

ORWC is worth its own subsection because it was implemented, iterated on for
roughly a year, and then deliberately removed. \commit{d565b2d} (2024-01-22) is the
first sign of trouble (\emph{``ORWC seem to have an issue''}); over the following
sixteen months, four separate fix attempts addressed a root-upgrade validation gap,
stripped redundant fences once OLC's own atomics were trusted, reworked guards from
raw pointers to cloned \path{SmartCell} handles after a dangling-pointer segfault on
free readers, and finally gave up on non-blocking semantics for the ordinary
mutable-borrow path, keeping only the optimistic \emph{upgrade} non-blocking.
\commit{3c6da7b} (2024-12-06) deprecates it; \commit{5c26abb} (2024-12-08) purges it
outright, $\sim$460 lines removed, with the commit message stating the reason
plainly: \emph{``duplicate with OLC.''} \textbf{Root cause, in one sentence:} ORWC's
optimistic upgrade path never reached a provably correct state, and once OLC was
fully validated it gave the identical guarantee through one simpler seqlock
mechanism --- ORWC became pure maintenance liability, not a genuinely distinct
protocol.

\subsection{The SmartCell latch}

\path{src/mv_sync/smart_cell.rs} implements the versioned latch every node uses: a
single 64-bit \path{cell_version} word (\path{OptCell<E>\{cell, cell_version\}}) with
three logical bit ranges packed into its top nibble --- a seqlock-style write flag,
pin-related flags, and a \emph{permanent} retired-tombstone bit. An optimistic read
(\path{borrow_read}) is one \path{Acquire} load with the write/retired bits masked
out, no lock taken. Upgrading to exclusive (\path{upgrade_write_lock}) is a CAS from
the captured read version to the write-flagged version, failing fast if the retired
bit is already set. Unlock stores \path{(write\_version+1) \^\ WRITE\_FLAG |
retired\_bit} with \path{Release}, deliberately reading the current retired bit with
a plain \path{Relaxed} load (safe, since nothing else can touch \path{cell_version}
while the writer guard is held) so a same-critical-section \path{mark_retired()}
call is preserved rather than clobbered by the unlock store (\S\ref{sec:gc} covers
the bug this specifically fixes). \path{try_retire} is a single CAS straight from a
reader snapshot to the retired state, skipping the intermediate exclusive step ---
sound only for guards that are unconditionally consumed with no back-out path.

\subsection{Chasing fences out of the hot path}

Once OLC's own atomics were trusted to carry the needed ordering, standalone memory
fences bolted on earlier, out of caution, were incrementally identified as redundant
and removed:

\begin{longtable}{@{}p{2cm}p{2cm}p{5.3cm}p{5.3cm}@{}}
\toprule
\textbf{Commit} & \textbf{Date} & \textbf{Change} & \textbf{Rationale} \\
\midrule
\endhead
\commit{48c4ded} & 2024-05-06 & Removed \path{fence(Acquire)}/\path{fence(SeqCst)} in
  the traversal loop & The surrounding OLC load/CAS already provides the ordering. \\
\commit{3c6da7b} & 2024-12-06 & Removed a paired fence duplicating
  \path{len.load(Acquire)} & An atomic load's own ordering already suffices. \\
\commit{a7bb29c} & 2025-04-01 & CAS failure ordering \path{Relaxed}$\to$\path{Acquire};
  unlock store \path{Relaxed}$\to$\path{Release} & Establishes the exact
  happens-before edges OLC readers rely on, cutting spurious CAS retries. \\
\commit{a518eed} & 2025-04-25 & Deleted standalone fences immediately adjacent to an
  atomic op on the same word & Pure dead weight once the adjacent op's ordering
  subsumed it. \\
\commit{2a2be93} & 2026-06-22 & Soft-commit length bumps \path{Release}$\to$\path{Relaxed}
  & The following write-unlock's own \path{Release} already publishes it. \\
\commit{40b00fd} & 2026-06-30 & \path{Node::m_type} \path{AtomicUsize}$\to$\path{usize};
  removed the \path{Drop} fence & Only mutated under the exclusive write lock; the
  latch's own release already publishes it. \\
\commit{45c7c21}/\commit{fd92f38} & 2026-07-21 & \path{try_retire} lets
  unconditionally-consumed guards skip a full write-lock upgrade during split/merge &
  Removes one CAS/lock-acquire per split/merge for guards with no backout path. \\
\bottomrule
\caption{Chronology of fence removal and ordering tightening.}
\label{tab:fences}
\end{longtable}

Alongside this, \commit{7f2bd09} (2026-05-21) cache-line-aligns \path{Node}
(\path{\#[repr(C, align(64))]}) so the hot, frequently-polled page-type tag does not
false-share with the rest of the node, and moves the root index onto the same
\path{SmartCell} machinery ordinary blocks use, off a separate \path{parking_lot}
mutex. A later experiment (\commit{45cb048}, 2026-05-28) that had range-query
iterators hold a long-lived \path{BlockGuard} across a whole iteration was
\emph{reverted}: a \path{'static}-lifetime guard held for an iterator's whole life is
exactly the kind of long-lived, never-revalidated read that GC's reclaim bookkeeping
does not know how to account for. The revert restores a bare \path{BlockRef},
re-borrowed fresh on each step, keeping GC's liveness accounting sound
(\S\ref{sec:gc}).

\section{Structural maintenance: splits, merges, and page capacity}
\label{sec:smo}

Structure-modification operations (SMOs, \path{src/mv_tree/smo.rs}) decide, on the
cold path, whether a page needs to split (overflow) or merge (underflow), and how.

\subsection{Capacity tuning}

\commit{34555e2} (2026-06-20) moved the internal-page overflow trigger from
\path{max_keys() - 3} to \path{max_keys() - 1}: a page is now allowed to fill much
closer to capacity before an SMO is attempted at all (\gainbounded{}). \commit{b29d347}
is the 4\,KiB-alignment trim already covered in \S\ref{sec:physical} (\gainmeasured{},
25\%$\to$100\% aligned allocations, $-9$--12\% dTLB load misses).

\subsection{Split: choosing where to cut}

\path{split()} decides between a \textbf{VERSION\_SPLIT} (compact dead versions out
of one page in place) and a \textbf{KEY\_SPLIT} (cut at a key boundary). Two real
bugs surfaced here under concurrent TPC-C stress in July--August 2026:

\begin{itemize}
  \item \commit{2a2607e} (2026-08-02) --- the VERSION\_SPLIT/KEY\_SPLIT decision used
        a strict \path{survivor\_count > capacity} test, so a page whose survivors
        landed at \emph{exactly} capacity wrongly took the ``just compact'' path,
        producing an already-100\%-full result with no room for the pending write ---
        a panic on a root-leaf tree, an infinite retry otherwise. Fixed with
        \path{>=} (\gainenabling{}). This turned out to be the true root cause of a
        stress-test failure originally suspected to be hot-key version tearing
        (\S\ref{sec:selfoverwrite}) --- the failing page in fact held several
        distinct keys.
  \item \commit{dc30111} (2026-08-02) --- the boundary search greedily combined
        ``don't tear a key's version chain'' and ``fit within capacity'' and could
        miss a valid boundary satisfying both. Fixed with a tiered search: (1) no
        tear and fits, (2) fits but may tear, (3) no tear but may overflow, (4) raw
        midpoint fallback (\gainenabling{}).
\end{itemize}

\subsection{Merge: choosing what to combine}

\begin{itemize}
  \item \commit{dcb5abe} --- compares available room on \emph{both} adjacent
        siblings instead of defaulting to the right one (\gainbounded{}).
  \item \commit{4c01360} --- dropped a \path{ts\_start}-based re-sort when combining
        two leaves: it risked reordering a live record ahead of a dead predecessor
        with a larger \path{ts\_start}, corrupting the backward-scan invariant from
        \S\ref{sec:physical} (\gainbounded{}, and a correctness fix).
  \item \commit{4bc78de} --- merging two leaves by \path{active\_count} alone could
        overflow \path{NUM\_RECORDS}, because \path{record\_survives\_gc} also keeps
        GC-protected dead records an active-count threshold cannot see. Fixed with an
        explicit \path{leaf\_merge\_would\_overflow} preflight that falls through to
        a key split when the true survivor count would not fit (\gainenabling{}).
  \item An experiment retrying an alternate merge candidate immediately
        (\commit{0659d34}) was implemented and \emph{reverted}: it spun without
        backoff and starved threads under contention.
\end{itemize}

\subsection{A week of hardening (2026-07-31 -- 08-02)}

The fixes above were the tail of a dense, one-week debugging arc chasing two real
TPC-C crashes: \commit{7432525} hardened \path{InternalPage} overflow checks without
yet fixing the crash; \commit{5c54658} fixed both via ``GC-gated survivor
protection'' (correctly distinguishing a genuinely committed delete from one whose
abort-reversal still needs its predecessor intact) and corrected \path{bulk\_push}
accounting (it had assumed every bulk-pushed record was active); \commit{77a5689}
reworked split/merge for a weakened fresh-node rule; \commit{fc8806b} reverted a
defensive \path{Result}-based workaround layered on \path{on\_overflow\_node}/
\path{on\_underflow\_node} once the real cause was actually found and fixed.

\section{The self-overwrite optimization}
\label{sec:selfoverwrite}

This is the second load-bearing optimization singled out in \S\ref{sec:foundational}:
\emph{when a single, still-uncommitted transaction updates the same key several
times, overwrite the previous update in place instead of appending a new
multiversion entry.}

\subsection{The problem}

Every write mints one more physical \path{RecordPoint} tuple, and an
\emph{uncommitted} tuple can be neither discarded (the transaction might still
commit) nor coalesced (nothing has superseded it, since nothing else can even see
it yet). The old update path was literally ``delete old, push new'' on every call,
so one open transaction writing the same key $N$ times left $N$ physical,
uncommitted, unreclaimable tuples for that one key. If $N$ grows large relative to a
page's capacity, \path{split()}'s boundary search has no key boundary to cut at ---
its own documentation states it falls back to tearing the chain across two
disjoint-fence siblings only when the page is entirely one key's chain, silently
orphaning the torn-off half from any fence-routed lookup. A concrete real-world
trigger: TPC-C's New-Order transaction pricing two order-lines for the same item
lands two updates on the same \texttt{Stock} row inside one transaction.

\subsection{The fix, as it exists in the current code}

The logic now lives in free functions shared by both \path{DbTransaction} and
\path{TpccTxn} (\path{src/mv_db/transaction.rs}, lines 92--216 as of this writing).
The update path:

\begin{lstlisting}
match leaf_page.latest_position(key, true) {
    Some(position) => {
        if tree.is_visible_stamp(worker_id, ts_start,
            leaf_page.version_at(position).insertion_stamp())
        {
            let stamp = TxStamp::new(worker_id, ts_start);
            tree.wal_log_write(stamp, |_| CRUDOperation::Update(key, payload.clone()));

            // Self-overwrite fast path -- see `insert_on_tree`'s doc.
            if leaf_page.version_at(position).insertion_stamp() == stamp {
                leaf_page.set_payload_at(position, payload);
                return (CRUDOperationResult::Updated(stamp.ts_start()), false);
            }

            // Otherwise: normal versioned update (delete old, push new).
            if !leaf_page.version_mut_at(position).delete(stamp) {
                return (CRUDOperationResult::ZeroAffected(KeyAlreadyDeleted), false);
            }
            let current_len = leaf_page.len();
            leaf_page.push_uncommitted(
                RecordPoint::new(key, VersionInfo::new(stamp), payload), current_len);
            leaf_page.commit_delta(0, 1);
            (CRUDOperationResult::Updated(stamp.ts_start()), true)
        } else { (CRUDOperationResult::Conflict, false) }
    }
    None => (CRUDOperationResult::ZeroAffected(KeyDoesNotExist), false),
}
\end{lstlisting}

\textbf{Detection} is a single equality check, not a heuristic: a \path{TxStamp}
packs \path{(worker\_id, ts\_start)} and is constant for the \emph{whole}
transaction, so ``was the newest live version of this key written by me, in this
very transaction'' reduces to \path{insertion\_stamp() == stamp} --- an OSIC concept,
as \S\ref{sec:foundational} points out. \textbf{Safety} follows from visibility:
\path{is_visible_stamp} only returns true for the same worker or a stamp already in
the commit log; neither applies to any other transaction while this one is open, so
there is no possible reader relying on this specific uncommitted tuple's current
content. If the found record belongs to another transaction, or is an
already-committed version, the code falls through to the ordinary
delete-then-push path unchanged.

The boolean second return value (\path{false} on the fast path) tells the caller
whether a new \path{written}-set entry is needed: mutating an already-open
transaction's own uncommitted record in place mints no new physical version, so the
entry an earlier write already pushed for this key still covers abort.

\subsection{Delete/reinsert cycles use the same bound}

\path{insert_on_tree} (lines 92--151) applies the symmetric case: the first
\path{delete(key)} followed by \path{insert(key, value)} in one transaction must
append a replacement tuple (the committed predecessor must stay physically present
for abort's undelete). But repeating that pair in the same transaction used to
append another replacement every cycle. The current code finds the newest
non-invalid record and, if its insertion \emph{and} deletion stamps both equal the
current transaction's stamp, clears the deletion mark and replaces the payload in
place instead of appending --- deliberately without adding another write-set entry,
since the first replacement's entry already makes abort invalidate that tuple and
undelete the original predecessor. $N$ delete/reinsert cycles now create one
replacement tuple; commit exposes only the final payload; abort restores the
pre-transaction value; WAL recovery replays the final logical state.

\subsection{The chronological story}

Two mechanisms are easy to conflate; only the second is the trick above.

\paragraph{An older, unrelated mechanism: GC-driven update-in-place (2025--2026).}
\commit{bea1117} (2025-06-05) added an update-in-place optimization to the
\emph{non-transactional}, single-operation CRUD dispatch path
(\path{src/mv_query/dispatch.rs}): if a record's insert version predates every live
snapshot --- GC has already proven nobody can need the old \emph{committed} value
--- its payload is overwritten in place rather than versioned again. It is gated off
whenever a WAL is attached (the shortcut cannot be represented as a single WAL
record) and remains in the codebase, but it is a GC/version-bloat heuristic for
single-op writes, unrelated to the transactional trick above.

\paragraph{The actual fix (August 2026).} \commit{160186e} (2026-08-01) first found
that \path{abort()} only reverted the \emph{newest} of several same-key writes,
leaving older self-written versions permanently visible; fixed to unwind the whole
chain one physical version at a time. \commit{85ac1b0} (2026-08-01) is the fix
itself, framed exactly via the TPC-C Stock-row scenario; its commit message also
documents a \emph{reverted} companion idea --- lowering the leaf overflow trigger to
buy compaction more headroom --- abandoned because it caused a hang under concurrent
stress. \commit{fbedd74} added a minimal repro isolating a \emph{different}
contributor to the same class of pressure (a long-open snapshot keeping an old
committed version GC-protected while other transactions keep committing new
versions of the same hot key --- no same-transaction repeats involved; the repro's
own documentation states self-overwrite ``cannot help here''). Tracing that repro
further (\commit{2a2607e}) found the real cause was the split off-by-one from
\S\ref{sec:smo}, not hot-key tearing at all.

\subsection{Resulting bound, and a known residual}

$N$ updates or delete/reinsert cycles to one key by one transaction now cost at most
one unresolved replacement tuple. First-writer-wins prevents another transaction
from concurrently extending that same unresolved chain. Once the writer resolves,
ordinary GC compaction discards obsolete committed versions; pages with multiple
keys remain splittable at a key boundary. The known, documented residual: writing
the same key $N$ times in one transaction still leaves \emph{one} extra physical
entry versus never writing it at all --- self-overwrite bounds the growth to $O(1)$,
it does not eliminate it. The TPC-C stress invariant still permits a narrow 0.5\%
difference between aggregate \path{Stock} counters and inserted \path{OrderLine}
rows, tracked as a separate, still-open accounting discrepancy unrelated to
same-key splitting or reader retention.

\section{Garbage collection: reclaiming dead blocks without contention}
\label{sec:gc}

\subsection{From a global mutex to sharded, lock-free reclaim}

\commit{acdc7a8} (2024-02-14) is the first GC implementation: a global
\path{Mutex<RBTree<SnapShot>>} of active transactions and a global
\path{Mutex<LinkedList<(*mut Version, BlockRef)>>} of dead pages, with the commit
message already naming the problem that would dominate the next several months:
\emph{``locks must be left out, or possible deadlocks appear!''} \commit{2e3ad21}
switched allocation to \path{try_lock()} on both structures (bail out rather than
block, \gainenabling{}); \commit{95e149f} moved the active-tx index to the project's
own OLC-locked concurrent B+-tree (itself later superseded). The current end state
(\path{src/mv_gc/block_tracer.rs}) drives liveness directly off
\path{TxContext::live_min_snapshot()} (\S\ref{sec:sharding}) rather than a separate
active-tx structure, and reclaim off a worker-sharded \path{SkipMap} rather than one
list or tree.

\subsection{Two ThreadSanitizer races, closed structurally}

\begin{itemize}
  \item \commit{d0e0977} (2026-07-21) --- a separate \path{retired: AtomicBool} field
        was folded into the SmartCell's own \path{cell_version} bits
        (\S\ref{sec:cc}). Before this, a two-step ``mark retired, then let
        \path{Drop} unlock later'' pattern raced a concurrent reuse's
        \path{clear_retired()} against the dying guard's deferred \path{Drop} storing
        a value computed at lock-\emph{acquisition} time, which could silently
        revert a freshly reused cell's version --- confirmed empirically as
        hangs/livelocks under GC plus heavy concurrency.
  \item \commit{ecc9900} (2026-07-27) --- two independent bugs. First, the dead-block
        index was keyed by a bare \path{Version}, but a merge's two dying siblings
        deliberately share one death version (the replacement block's birth
        version); \path{SkipMap::insert} on a duplicate key silently displaces the
        existing entry, leaking one of the two retired blocks. Fixed by keying on
        \path{(Version, block\_address)}. Second, \path{insert\_stamp}/
        \path{delete\_stamp} were plain, non-atomic fields mutated by
        \path{delete}/\path{invalidate} while lock-free OLC readers read the same
        word with no lock --- a genuine torn-read race, confirmed by
        ThreadSanitizer. Fixed by backing both with \path{AtomicU64}/\path{Relaxed}
        (sufficient because the hazard was torn reads, not ordering).
  \item \commit{fe03d90} (2026-07-28) --- the capstone: \path{version\_region} moves
        from \path{AtomicVersion} to plain \path{Version}, and the in-place
        \path{mark\_version\_obsolete} \path{fetch\_or} (itself confirmed racing
        range-query iteration) is deleted outright, replaced by the
        \path{live\_mask()} derivation described in \S\ref{sec:physical}. No writer
        ever mutates a slot after publication anymore, so the race does not merely
        get fenced correctly --- it stops being possible.
\end{itemize}

\textbf{Recurring pattern.} Nearly every ``racey''/``stabilize'' fix in this section
follows the same shape: something was mutated \emph{in place} on an already-published,
lock-free-readable slot, requiring ever more careful atomic/fence discipline to patch
around it, until the design changed so nothing is mutated after publication at all.
This mirrors \S\ref{sec:cc}'s fence-removal trajectory: converging from ``atomics and
fences bolted on to patch races'' toward ``structure the data so the race is
structurally impossible.''

\subsection{Sharded reclaim, historical routing, and abort-safe retention}

\path{block_tracer.rs} holds one \path{SkipMap} shard per worker
(\S\ref{sec:sharding} covers the mechanism); a nonblocking reuse fallback allocates
fresh memory rather than joining the page-latch dependency graph when reuse metadata
is contended. Old readers are kept on their retired \emph{source} block via
historical root/version routing rather than having every old-reader-visible record
copied forward into replacement pages during a rebuild --- copying for the globally
oldest reader would couple every key's reclaim to one transaction and risks a
self-reinforcing compaction livelock (a version of exactly this failure mode is
discussed as an open issue in \S\ref{sec:open}).

\section{Root indexing strategies}
\label{sec:root}

A reader must find ``the root as of snapshot $S$'' before traversing anything.
\path{RootIndexType} selects one of four version-keyed structures:

\begin{center}
\begin{tabular}{@{}p{2.6cm}p{4.2cm}p{2.9cm}p{4.0cm}@{}}
\toprule
\textbf{Strategy} & \textbf{Structure} & \textbf{Lookup} & \textbf{Notes} \\
\midrule
Vanilla & \path{LinkedList<Root<..>>} & $O(n)$ linear & Simplest baseline. \\
SK (SkipList) & \path{crossbeam\_skiplist::SkipMap} & $O(\log n)$ expected &
  Lock-free concurrent inserts/lookups. \\
DexaBTree & external OLC-locked B+-tree, keyed by \path{SnapShot} & $O(\log n)$ &
  Reuses a full concurrent B+-tree just to index root pointers. \\
FrugalList (default) & lock-free singly-linked list with a probabilistic
  skip-tower ($p=0.5$) & sub-linear expected & Appends are one \path{AtomicPtr}
  store (no CAS loop needed --- already serialized by the enclosing write latch);
  much smaller per-node footprint. \\
\bottomrule
\end{tabular}
\end{center}

These are not partitioned by version range --- each is one tree-wide structure
holding every historical root. They optimize a different axis than
\S\ref{sec:sharding}: the cost of finding the root pointer itself (a potential
single hot spot every reader must consult), trading lookup complexity against
per-append cost and memory overhead.

\section{Sharding as a system-wide pattern}
\label{sec:sharding}

The word ``sharding'' appears in four genuinely distinct places, none of them
classic horizontal key-range partitioning of one index. Each instead partitions a
specific hot, shared piece of \emph{bookkeeping state} by worker (or, in one case,
by table). \S\ref{sec:foundational} argued that this pattern presupposes both OLC
(so there is lock-free contention worth removing) and a stable per-worker identity
(so there is something to shard along); this section is the pattern applied four
times.

\subsection{Tables as shards: one tree per table}

\commit{5f92e9f} (2026-07-15) --- before this, multiple logical tables were
multiplexed into a single tree. \path{Database} now holds
\path{tables: ArcSwap<TableList<..>>}: each table gets its own tree, root index, and
--- importantly for \S\ref{sec:gc} --- its own dead-block tracker. Concurrent
transactions touching different tables no longer contend on one tree's root/internal
latches or one dead-block skip list. What \emph{must} stay global for cross-table
snapshot isolation --- the clock, commit-log pruning, snapshot liveness --- moves
into one shared, non-generic \path{TxContext} introduced in the same commit
(\gainbounded{}).

\subsection{GC dead-block tracker: own-shard-first, round-robin steal}

\path{BlockTrace} holds one \path{SkipMap} shard per CPU
(\path{shard\_for(worker\_id) = worker\_id \% shards.len()}). Registering a died
page only ever inserts into the dying worker's own shard; reclaim pops from the
caller's own shard first, falling back to a shared round-robin cursor across other
shards only on a miss. \textbf{Why:} death versions all come from one global clock,
so before sharding, concurrent SMOs from every worker inserted into a \emph{single}
skip list clustering near its current maximum key --- an always-contended tail.
Since a reclaimer only needs \emph{a} reclaimable block, not the globally oldest
one, trading global pop-order precision for per-worker insert locality is sound.
\commit{f9fc960} (2026-07-23) added the ``try own shard first'' fast path (before it,
reclaim always scanned round-robin even when the caller's own shard already had a
usable entry).

\subsection{\texttt{TxContext}: cache-padded, per-worker atomic slots}

Three parallel vectors, one slot per \path{WorkerId}: \path{live_tx} (each worker's
currently published outermost snapshot), \path{live_tx_depth} (reentrancy depth, so
only the outermost \path{begin_snapshot} touches \path{live_tx}), and
\path{in_flight_bound} (a conservative lower bound published while a worker is
mid-registration --- has drawn a \path{ts\_start} but not yet inserted). Each worker
only ever writes its own slot; \path{free_block}/\path{live_min_snapshot} scan all
slots with \path{Acquire} loads, $O(\mathrm{max\_workers})$, no cross-slot
synchronization. This directly replaced two shared, globally-contended structures,
and the commit history records the profiling numbers that motivated it: \path{live_tx}
replaced a shared \path{SkipMap<SnapShot, AtomicUsize>} measured (per an in-code
comment) eating ``roughly a quarter to half of all CPU cycles'' on a saturated TPC-C
run; \path{in_flight_bound} replaced a single global \path{AtomicUsize} that, under a
16-thread update-heavy YCSB workload, rejected 73.8\% of \path{free_block} attempts
(versus 5.2\% at 2 threads). \commit{eac6e40} (2026-08-02) then wrapped all three
vectors in \path{crossbeam\_utils::CachePadded}: previously up to 8 adjacent
workers' slots could share one 64-byte cache line, directly contradicting those
fields' own design-intent comments --- a pure layout fix targeted at
64-core/128-thread scale, where cross-CCD coherence traffic is expensive.

\subsection{Thread-local worker identity and snapshot cache}

\path{src/mv_sync/worker.rs} hands out a \path{WorkerId} once per (tree, thread),
cached thread-locally keyed by a process-wide \path{tree\_uid}. A second thread-local,
\path{SNAPSHOT\_CACHE}, gives each worker its own \path{SnapshotCache} --- the OSIC
paper's ``thread-local snapshot cache'' --- rather than a \path{Vec} on the tree
indexed by \path{WorkerId}: using actual thread-local storage means no thread ever
\emph{could} touch another worker's cache, stronger than ``happens to always pass
its own id.'' A dedicated regression test file (\path{tests/todays\_optimization\_
regression\_tests.rs}) verifies, among other things, that a thread cycling through
three exactly-sized tree registries never loses a \path{WorkerId} across the two
inline TLS slots' lazy overflow path.

\section{OSIC integration: instant commit, lazily hardened}
\label{sec:osic}

\subsection{What OSIC is, and how cMVBT realizes it}

Ordered Snapshot Instant Commit~\cite{alhomssi2023scalable} decouples visibility
from durability. Every transaction draws two timestamps from one global,
monotonically increasing counter (\path{GlobalClock}, a single
\path{AtomicVersion}, \path{fetch\_add(1, SeqCst)}): \path{ts\_start} at begin,
\path{ts\_commit} at commit. Each worker keeps its own append-only \path{CommitLog}
(\path{Mutex<Vec<Version>>}); \path{commit()} draws \path{ts\_commit} then pushes it
under the lock --- that ordering is exactly what prevents a reader from observing a
reserved-but-not-yet-inserted stamp. Visibility (\path{is_visible}) rejects invalid
(aborted) stamps outright; same-worker writes are visible if
\path{stamp.ts\_start() <= reader\_ts\_start}; otherwise visibility is
\path{LCB(stamp.worker\_id, reader\_ts\_start) > stamp.ts\_start()}, computed via
each thread's own \path{SnapshotCache} memo. \path{WalWriter} is a lock-free MPSC
channel to a background flush thread: \path{log\_with\_stamp}/\path{log\_commit}
encode and enqueue, returning immediately --- fire-and-forget, no per-op fsync wait.
\path{flush\_loop} sleeps a short \path{GROUP\_COMMIT\_LINGER} ($200\,\mu$s) after
the first message, drains everything else non-blockingly, batches it into one
\path{write\_all} + \path{sync\_data()}, and advances a \path{hardened} watermark to
the maximum \path{ts\_start} of any \emph{Commit} message in that batch. Recovery
keeps only writes with a matching commit marker (commit-gated replay), sorted by
\path{(ts\_commit, file order)} since commit order can differ from log order across
concurrent workers. \path{DbTransaction::commit} appends to the committing worker's
\path{CommitLog} --- making every write visible everywhere, instantly --- then logs
exactly one fire-and-forget WAL commit marker. \emph{This is the whole ``instant
commit'':} one in-memory append, with nothing on the write/commit path waiting on
disk I/O; durability is a separate, batched, background concern, exposed only via an
explicit \path{wait\_wal\_hardened()} API for callers who specifically need a
durability checkpoint.

\paragraph{Evolution.} A pre-OSIC era (Dec 2025 -- Jan 2026) used a hand-rolled
sharded thread-indexed ``committed'' watermark array --- in effect an imperfect
reinvention of the CommitLog/GLC pair, with a sequence of patches fixing edge cases
around threads changing roles or exiting, reading as a slow discovery, by trial and
error, of exactly the problem OSIC's design solves cleanly. \commit{714a82f}
(2026-07-10) replaced it wholesale with the paper's actual design.

\subsection{The WAL performance campaign}

A short-test profiling pass (\path{scripts/compare_engines.py}, \path{perf record}/
\path{perf annotate} on a debug-symbol release build) found three concrete issues,
fixed together and re-measured on the same harness invocation.

\textbf{1. \path{target-cpu=native} was silently dead.} \path{Cargo.toml} had a
\path{[build] rustflags} block, but Cargo only reads \path{rustflags} from
\path{.cargo/config.toml}, never the manifest --- the whole binary had been
compiling to the generic x86-64 baseline. Fixed by moving the flag into a new
\path{.cargo/config.toml} (\gainenabling{}, confirmed present in the current
repository).

\textbf{2. Hand-rolled CRC32 was the single biggest hotspot.} A byte-at-a-time,
8-shift-per-byte bit loop was \textbf{22.6\% of all CPU cycles} (self time) on
YCSB-A at 64 threads. Replaced with a 256-entry lookup-table CRC32, same
polynomial/format, pinned against the standard \path{"123456789"}
$\to$ \path{0xCBF43926} check value. Self-time dropped to 11.3\% on the same
workload.

\textbf{3. The single shared WAL channel was a real contention point.}
\path{Sender::send} showed 86\% of its own samples in a \path{crossbeam-channel}
internal spin-wait, worth $\sim$8\% of \emph{all} CPU cycles on TPC-C at 64 threads.
Fixed by splitting \path{WalWriter} into \path{NUM\_SHARDS = 4} independent
(channel, flush-thread) pairs keyed by \path{worker\_id \% 4}, all still writing
through one shared \path{Arc<Mutex<File>>} so on-disk format and recovery are
unaffected; the shared \path{hardened} watermark switched from \path{store} to
\path{fetch\_max} since multiple flush threads can now finish out of order.

\begin{center}
\begin{tabular}{lrrr}
\toprule
Workload & Threads & Before & After (\gainmeasured{}) \\
\midrule
tpcc & 16 & 40,902 & 43,913 (+7.4\%) \\
tpcc & 64 & 57,667 & 67,812 (\textbf{+17.6\%}) \\
ycsb\_a (50\% update, Zipfian) & 16 & 1,603,135 & 2,010,940 (+25.4\%) \\
ycsb\_a & 64 & 1,339,295 & 2,258,954 (\textbf{+68.7\%}) \\
ycsb\_e (scan-heavy) & 16 & 2,137,315 & 2,483,067 (+16.2\%) \\
ycsb\_e & 64 & 1,789,143 & 2,679,747 (\textbf{+49.8\%}) \\
ycsb\_f (read-modify-write) & 64 & 1,355,666 & 2,305,318 (\textbf{+70.1\%}) \\
\bottomrule
\end{tabular}
\end{center}

The read-only/read-mostly workloads (ycsb\_b, ycsb\_c) moved by under 1.2\% either
way, confirming the fixes are surgical to the write path. The clearest structural
result: the two workloads that previously got \emph{worse} with more concurrency
(ycsb\_a, ycsb\_e) now scale up instead of down. Not addressed by this pass: TPC-C's
OCC conflict rate remained very high at 8 warehouses/64 terminals (Payment
$\sim$71--76\% of attempts conflict and are discarded; New-Order $\sim$47\%) ---
flagged in \S\ref{sec:open}.

\subsection{A lock-free WAL backend, and two concurrency bugs it took to land}

A second, concurrent line of work attacked the same contention differently:
\path{LockFreeWalWriter}/\path{LockFreeWalBackend} removes the channel and
dedicated writer thread entirely --- every worker reserves its own byte range via
\path{tail.fetch\_add} and calls \path{pwrite} directly, optionally batching its own
records locally (\path{LocalBatch}) before flushing. \path{WalBackend} wraps both
designs behind one enum so callers pick either at attach time.

Merging this (\commit{ebdd789}) with the sharded-channel work (\commit{4a1a50a})
surfaced two bugs unrelated to either design's core idea:

\begin{itemize}
  \item \textbf{An OOM, not a WAL bug.} \path{mv\_test::RESTART\_TRACE} (a
        write-restart diagnostic, documented ``off by default'') had been left
        \path{true} in checked-in code. Every OCC restart recorded a \path{String}
        key into a process-lifetime global map that nothing ever clears; a
        comparison loop calling \path{run\_tpcc()} six times in one process
        accumulated across every iteration until the fastest config grew past 12\,GB
        RSS and got OOM-killed, taking the terminal session with it. Fixed by
        restoring the documented default and adding \path{reset\_restart\_trace()},
        called at the start of every \path{run\_tpcc()} (verified present in the
        current source, \S\ref{sec:fresh}).
  \item \textbf{A cross-shard \path{hardened\_version} race.} The sharded
        \path{WalWriter} published every shard's own confirmed watermark into one
        value via \path{fetch\_max} --- correct only if every shard's watermark were
        directly comparable, which it is not, since \path{ts\_start} is one global
        sequence spread across shards unpredictably. One busy shard's flush could
        advance the shared watermark past a \path{ts\_start} a slower shard had not
        flushed yet; a consistency test caught it directly, a quarter of written
        records missing from the reconstructed WAL. A first fix attempt (minimum of
        every shard's raw confirmed value) traded the false-positive for a hang: a
        shard whose submissions never reach some other shard's higher watermark
        would never satisfy a plain minimum, even once fully caught up. The actual
        fix (\commit{9ced1a3}) gives each shard a \path{submitted} watermark
        alongside \path{confirmed}, and treats a shard as imposing no constraint
        once \path{confirmed >= submitted} (fully drained) rather than comparing raw
        values across shards. Verified: the failing test 5$\times$ clean, the full
        suite 3$\times$ clean.
\end{itemize}

\textbf{Synthetic backend comparison} (original report, debug build):

\begin{center}
\begin{tabular}{lrr}
\toprule
Backend & @8 threads & @16 threads \\
\midrule
\path{WalWriter} (4-way sharded) & 1,006,755 ops/s & 1,353,816 ops/s \\
\path{LockFreeWalWriter}, unbatched & 752,179 ops/s & 748,676 ops/s \\
\path{LockFreeWalWriter} + \path{LocalBatch(64)} & 2,549,030 ops/s & \textbf{5,089,502 ops/s} \\
\bottomrule
\end{tabular}
\end{center}

Unbatched lock-free writes lose to the sharded channel past a few threads: each
write pays a full \path{pwrite} syscall plus filesystem-level inode-extension
serialization on a shared growing file. Per-thread batching is what actually wins.
\S\ref{sec:fresh} reproduces this exact comparison, in release mode with LTO, with a
fuller thread/batch sweep, on the machine this document was written on.

\subsection{Ticket-less logging and exact framing hints}

A follow-up \path{perf record} pass, once the channel contention above was no longer
dominant, found two purely allocator-side hotspots it had been masking:
\path{WalWriter::enqueue} built a fresh \path{bounded(1)} ack channel on
\emph{every} call to return a flush ticket that every production call site always
discarded; and every framing buffer's capacity hint (24--44 bytes) was tuned for the
module's own \path{u64}-payload tests, while \path{YcsbRow} is roughly 25$\times$
that and \path{TpccRow} variants routinely 7--10$\times$ over, forcing several
grow-and-copy reallocations per WAL write. Fixed by splitting logging methods into a
ticket-less default plus an explicit \path{*\_with\_ticket} variant, and by adding
\path{WalPayload::wal\_encode\_size\_hint()} with exact overrides mirroring each
payload type's own encoder.

\begin{center}
\begin{tabular}{lrrr}
\toprule
Metric & Before & After & Change \\
\midrule
TPC-C, WAL on (tpmC) & 2,675,157 & 3,013,193 & \textbf{+12.6\%} \\
YCSB-A, WAL on (ops/sec) & 1,374,921 & 1,723,455 & \textbf{+25.3\%} \\
Synthetic \path{WalWriter}, 16 threads, \path{u64} payload & 1,353,816 ops/s & 2,427,071 ops/s & \textbf{+79\%} \\
\bottomrule
\end{tabular}
\end{center}

The synthetic case gained the most because, against an 8-byte payload, the discarded
ticket's channel allocation was close to the entire per-write cost. \textbf{Not
fully closed:} a fresh $\sim$1000-byte \path{Vec} is still allocated per WAL write on
\path{WalWriter}'s cross-thread path; a true buffer pool was flagged as a follow-up
and remains unimplemented (\S\ref{sec:open}).

\subsection{Two dead ends, honestly reported}

\textbf{A thread-local scratch buffer on \path{LockFreeWalWriter}} was implemented
and verified correct, on the reasoning that the lock-free writer's synchronous
\path{pwrite} has no cross-thread ownership problem. Tracing the actual call graph
found it was dead weight: every real caller goes through \path{LockFreeWalBackend},
which always routes into \path{LocalBatch}'s own persistent buffer first, even at
\path{batch\_size: 1} --- so only this module's own unit tests and the deliberately
unbatched benchmark arm ever reached the patched method. \textbf{Reverted} rather
than kept with no real effect.

\textbf{An opt-in \path{mimalloc} allocator feature} targeted the genuine
cross-thread free pattern in \path{WalWriter} (one thread mallocs the framed
buffer, the flush thread frees it). Measured end to end: YCSB-A gained 2.3--3.0\% on
the one backend with that cross-thread free, a wash on backends without it; TPC-C
\emph{lost} 3.3--4.9\% across every backend, including ones with nothing to do with
this allocation pattern, all six (backend $\times$ terminal-count) combinations
moving the same direction. \textbf{Not adopted as default} (jemalloc stays the
global allocator); kept as an opt-in feature since the trade-off is genuinely
workload-dependent, not a universal win.

\subsection{Atomic operations: an auto-commit fast path alongside real transactions}

A single-key \emph{atomic} operation is auto-committed, not a snapshot transaction
with a one-entry write set. Its write path: take a short reclamation pin during
traversal, acquire the target leaf's exclusive OLC writer guard, verify the newest
physical version is visible at the current OSIC boundary, draw the stamp and log the
WAL write, mutate the leaf while unobservable behind the guard, append to the
worker's commit log, \emph{then} release the guard --- publishing an
already-committed version. The leaf-guard release is the linearization point for
other OLC traversals; by the first instant a reader can validate the modified leaf,
the new record's stamp is already in the commit log, closing the ``published leaf
entry, commit log not updated yet'' window a naive design could otherwise expose.
Deliberately, an atomic write does \emph{not} call \path{begin_snapshot}: it wants
the newest committed state rather than a historical view spanning several
statements, and it retains no root/page/timestamp after the call, so registering a
full long-lived snapshot would add clock-tick and bookkeeping cost the protocol does
not need. The YCSB driver exposes both modes (\path{--cmvbt-ycsb-mode atomic|
transaction}) without changing record shape, key distribution, or WAL settings,
specifically so the cost of the transaction machinery itself is measurable rather
than baked into every number.

\subsection{Known, documented limitations}

Unbounded \path{CommitLog} growth under long GC-off historic-query runs (a
chunked/segmented log was considered, explicitly deferred). \path{MAX_WORKERS_CAP =
256}, a generous but hard ceiling. Abort must revert writes in strict LIFO order ---
a documented but delicate invariant, not a generically solved problem. The
self-overwrite residual noted in \S\ref{sec:selfoverwrite}. And \commit{ecc9900}'s
torn-read fix is explicit in its own documentation about closing one specific
hazard, not every conceivable concurrent-mutation hazard on the lock-free OLC read
path.

\section{Range scans and analytical terminals}
\label{sec:scans}

\subsection{Why internal pages cannot be scanned in physical order}

An internal page is append/version ordered, not key ordered: structural
modifications append replacement child entries while retaining historical entries
older snapshots still need. Neither walking physical slots left to right nor sorting
every snapshot-matched slot produces a correct scan --- the latter includes
superseded, overlapping children. \path{RangeQueryIter} instead retains internal
pages on its current path: at each level it searches entries in reverse append
order and selects the newest snapshot-visible fence containing the scan cursor,
descends, and after consuming a leaf advances the cursor one key past that leaf's
fence and reroutes from the retained parent --- the same selection rule as point
lookup, allocating no child lists. An exact full-domain scan
(\path{Key::MIN..=Key::MAX}) is detected once before the leaf loop and uses a
visibility-only record loop, skipping two bound comparisons per physical record that
every bounded interval still needs. The final leaf is a special case: key increment
saturates at the maximum key, so both the materialized and streaming paths must
explicitly complete when the consumed fence reaches the requested upper bound or the
tree maximum, or a full-range streaming scan would route the unchanged maximum
cursor back to the final child forever.

\subsection{Rejected: child-list reuse, small-page dedup, hybrid traversal, lean frames}

Four related optimizations were implemented, benchmarked (an experimental version
appeared 18--20\% faster), and then \textbf{removed} once GC stress testing proved
their shared premise incorrect: reusing an eagerly constructed child-list snapshot,
deduplicating boundaries by hashing/linear scan, a hybrid narrow-range traversal
built on that same list, and releasing an internal page early after expansion.
Equal-boundary deduplication cannot prove that arbitrary historical intervals are
fully superseded; under copy-on-write GC it let both a current route and an
overlapping historical route be scanned, returning duplicate logical keys. The
active cursor-based traversal above was already allocation-free and always selects
the newest visible route, making the rejected machinery unnecessary as well as
unsound.

\subsection{Streaming terminals: zero-copy analytical consumption}

\path{RangeQueryIter::for_each_ref}/\path{fold_ref}/\path{try_for_each_ref}/
\path{count_ref} call a visitor with \path{(Key, \&Payload)} directly, avoiding one
\path{RecordPointResult} construction, one payload-handle clone, a retained leaf
buffer, and materializing a whole result vector, per match. The payload reference is
callback-scoped and must not be retained. \path{DbTransaction} exposes the same four
operations (\path{range\_count}, \path{range\_fold}, \path{range\_for\_each},
\path{try\_range\_for\_each}), keeping its function-local table \path{Arc} alive
until the terminal operation consumes the scan; \path{TpccTxn}'s versions dispatch
across both standard and big-tree table classes (\S\ref{sec:bigtree}) with identical
behavior.

\textbf{Measured} (200,000 shuffled \path{u64} rows, five warmups, 40 alternating
samples): streaming was 1.165$\times$--1.173$\times$ faster than the materialized
iterator for a full scan ($\sim$16.8\%; 28.8--29.3\,M rows/s vs.\ 24.7--25.1\,M
rows/s). A later, separately measured \path{count\_ref} build reached
\textbf{49--53\,M rows/s} on full scans and showed 1.18$\times$--1.37$\times$ speedups
on both full and narrow-range scans; these routing-adjacent numbers are retained
only as design history, since the specific implementation measured there (inline
reusable scratch, a single-child fast path) was the one later reverted for
correctness in the previous subsection. \path{count\_ref} itself remains enabled and
its own number is not affected by that revert.

\subsection{Fix 1: test the range before the visibility check}

\path{RangeQueryIter::refill} and \path{key_range_read_from_root} both filtered each
physical record with \path{r.version().matches(is\_visible) \&\& range.contains(r.key())}.
\path{VersionInfo::matches} calls \path{is\_visible} up to twice per record (insertion
and deletion stamp), on \emph{every} dead/superseded record a leaf holds --- the
majority of a garbage-heavy leaf's contents, per \S\ref{sec:bigtree}'s scan-cost
findings --- and it did this \emph{before} the cheap, inlinable range comparison
ever got a chance to reject the record. Swapped the order in both places: pure
reorder of an existing \path{\&\&}, zero behavior change. No benefit for a full-table
scan (its range always contains every key). Real benefit for a \emph{narrow} scan
sharing a leaf with keys outside its range --- e.g.\ CH-benCHmark Q4/Q5's per-order
\texttt{OrderLine} sub-scans. \textbf{Measured} (8 warehouses/16 terminals, 20s runs,
noisy/contended, indicative not averaged): Q4 latency $\sim$614--617\,ms
$\to$ $\sim$505--507\,ms (\textbf{$\sim$18\% faster}); Q5 $\sim$6.81--6.90\,s
$\to$ $\sim$6.09\,s (\textbf{$\sim$11--13\% faster}). Q1/Q6 (full-table scans) showed
comparable-or-noisy numbers either way, exactly matching the mechanism's prediction.

\subsection{Fix 2: remove the \code{dyn FnMut} indirection}

Reordering only helps \emph{out-of-range} records; every record actually in range
still paid for a non-inlinable indirect call on every \path{matches} invocation.
\path{VersionInfo::matches} was generalized from \path{\&mut dyn FnMut(TxStamp) ->
bool} to \path{<F: FnMut(TxStamp) -> bool + ?Sized>(\&mut F)} --- backward
compatible, since \path{dyn FnMut(..)} itself satisfies the new bound, so every
existing caller keeps compiling with zero behavior change. \path{RangeQueryIter::refill}
was restructured to call a new \path{with\_snapshot\_cache\_and\_logs} accessor once
per \path{refill()} (not once per leaf), build a concrete \path{is\_visible} closure
inline, and use it directly --- no \path{dyn} anywhere in this path. \path{query.rs}'s
two call sites were left completely untouched.

\textbf{A methodological scare, worth recording as a case study.} The first
verification pass looked like a real regression:
\path{verify\_concurrent\_shared\_keys}'s ``never loses an update'' stress test
flaked 2/5 then 4/20 with the change, against a first baseline sample of a clean
10/10. That pattern --- a concurrency test flaking only once concurrency-adjacent
code changes --- is exactly what a real bug from this kind of change would look
like. Re-investigating rather than reverting on the spot: a \emph{wider} baseline
sample, drawn after the same session's own repeated heavy test/build load had warmed
the machine up, came back 2/19 ($\sim$10.5\%) on the identical test and panic line
--- the earlier clean 10/10 baseline had simply been drawn before the machine warmed
up, not because the flake was not there. A clean, side-by-side 20-run batch with the
change applied came back 2/20 (10\%), statistically indistinguishable from baseline.
\textbf{Lesson:} do not trust a 5--10 run sample of a timed concurrency test, in
either direction, on a machine that has been running heavy benchmarks for a while
--- draw a same-conditions baseline of comparable size before attributing a failure
to a change. (This exact flake, and its true underlying cause, are followed further
in \S\ref{sec:open}.)

\subsection{The point-read miss bug this work incidentally surfaced}

Investigating the flake above led to a second, independent, and more serious bug,
now fixed and committed (\commit{c0fe172}): \path{key_point_read_from_root}'s
\path{skip_while} and \path{key_range_read_from_root}'s binary-search cutoff both
assumed a leaf's physical append order tracks \path{ts_start} order. That is false
under this system's concurrency model --- \path{ts_start} is drawn at \path{begin()},
\emph{before} the writer acquires the leaf's write lock to physically append, so two
concurrent writers can draw \path{ts\_start} in one order but append in the other.
When they diverged, the point-read shortcut could walk straight past the one visible
record and return an empty result for a live, never-deleted key. \textbf{Fix:}
removed both shortcuts; always scan the whole leaf --- the same cost
\path{RangeQueryIter::refill} already always pays, so no new asymptotic cost class
is introduced (\gainenabling{}). Verified with 20+20 clean runs of both affected
tests, versus the $\sim$10\% baseline failure rate documented beforehand.

\section{Workload and benchmark-harness optimizations}
\label{sec:workload}

Measurements in this section used a release build, 200,000 preloaded rows, four
workers, a three-second timed phase, Zipf $\theta=0.99$, GC enabled, no WAL, and the
frugal root index --- they measure a combined batch, not isolated attribution to one
row.

\begin{center}
\begin{tabular}{p{4.3cm}p{7.4cm}p{2.4cm}}
\toprule
\textbf{Change} & \textbf{Implementation and gain} & \textbf{Evidence} \\
\midrule
Bulk-generate YCSB values & Rows allocate their final payload once, filled in place
  from buffered random bytes. Full-row YCSB-A: 2.52M $\to$ 3.11M ops/s
  (\textbf{+23.5\%}). & \gainmeasured \\
\path{writeallfields} (standard mode) & Standard YCSB updates one field via a
  leaf-latched \path{update\_with} patch, not a full-row regenerate. Standard
  single-field YCSB-A: \textbf{4.81M ops/s} (+90.8\% vs.\ the old full-row baseline
  --- different semantics, not a fair like-for-like number). & \gainmeasured \\
\path{point\_exists\_si} & Hit/miss-only reads test visibility in place, no result
  vector or payload clone. Combined with alias Zipf sampling, YCSB-C: 7.16M
  $\to$ 14.82M ops/s (\textbf{+106.8\%}). & \gainmeasured \\
Walker-alias Zipf sampling & A finite alias table built once at setup; sampling is
  two random draws and a lookup, replacing per-request floating-point power
  computation. $O(n)$ setup, 8 bytes/key ($\sim$160\,MB at 20M keys). & part of above \\
Empty read-only commit elision & A read-only transaction with an empty write set
  releases its snapshot without drawing \path{ts\_commit} or appending a commit-log
  entry --- one clock tick and one commit-log op removed per such transaction. &
  \gainbounded \\
Sampled YCSB-E scan latency & One in 1,024 scans receives a latency timestamp; the
  timeseries clock refreshes every 256 ops. YCSB-E: 4.59M $\to$ 6.22M ops/s
  (\textbf{+35.4\%}). & \gainmeasured \\
Full-domain scan specialization & \path{try\_for\_each\_ref} detects
  \path{Key::MIN..=Key::MAX} once, uses a visibility-only leaf loop (see
  \S\ref{sec:scans}). & included above \\
Isolate CH Q1/Q6 & \path{OlapMode::ChQ1}/\path{ChQ6} run only the named query; the
  old rotating \path{ch} mode mixed in unrelated Q4/Q5 joins. & \gainenabling \\
\bottomrule
\end{tabular}
\end{center}

Also unrelated to raw throughput but worth recording: Zipfian rank scrambling
(FNV-based) was added specifically to make the benchmark's key-access locality
\emph{more} representative of real Zipfian workloads, not to make the engine
faster --- and \path{fastrand} replaced the earlier RNG for TPC-C parameter
generation purely to take generation cost off the transaction hot path. Both changed
the generated distribution, so before/after engine comparisons using either need to
pin the generator version. Finally, the \path{RESTART\_TRACE} OOM story from
\S\ref{sec:osic} directly motivated two always-on (not \path{\#[ignore]}d),
deliberately tiny WAL perf smoke tests, and the default TPC-C/YCSB comparison scale
was reduced to a few seconds per point specifically so the full suite (151 tests as
of this writing, \S\ref{sec:fresh}) runs on an ordinary development machine in
single-digit seconds; \path{CMVBT\_FULL\_BENCH=1} restores the original large-scale
configuration for dedicated performance runs.

\section{Case study: \code{BigTreeSize} --- a contention/garbage continuum, worked end to end}
\label{sec:bigtree}

This section is presented as a standalone case study because it is the clearest
example in the whole codebase of an optimization whose two extremes are each bounded
by a \emph{different, independently mechanistic} cost, and because the investigation
of it produced and then caught its own benchmark methodology bug --- a useful
lesson in its own right.

\subsection{Motivation}

Per-table root-restart attribution (\path{mv\_test::dump\_root\_restarts\_by\_table})
found that \texttt{district} and \texttt{warehouse} alone accounted for
\textbf{$\sim$99.97\% of all root restarts} at default TPC-C scale, despite being the
smallest tables in the schema (10 and 1 row per warehouse, respectively). The reason:
leaf overflow triggers at a hard 100\%-full ceiling with no tunable slack, and these
two tables are hit by write volume disproportionate to their row count, so their
physical version-chain growth crosses that ceiling roughly every 40 writes, forcing a
compaction that write-locks the one page nearly every concurrent transaction routes
through.

\subsection{The current design: an eight-variant size-class enum}

\path{BigTreeSize} (\path{src/mv_bench/tpcc_schema.rs}) gives \path{Table::Warehouse}/
\path{Table::District} their own tree type (\path{TreeClass}, plain tag dispatch, no
\path{dyn Trait}), independent of the base tree's own untouched 123-record leaf. As
of the current code, this is an eight-way, not five-way, continuum, spanning both
below and above the base tree's own size, each variant's \path{NUM\_RECORDS} chosen
so \path{OptCell<Block<..>>} lands exactly on a page-size multiple with zero waste:

\begin{center}
\begin{tabular}{lrrp{6.2cm}}
\toprule
Variant & Records/leaf & Page size & Measured trade-off (in-code, doc-commented) \\
\midrule
\path{KiB1} & 27 & 1\,KiB & Smallest variant confirmed to hold up under concurrent
  load --- see below for why nothing smaller exists. \\
\path{KiB2} & 59 & 2\,KiB & \\
\path{KiB4} & 123 & 4\,KiB & Same shape as the untouched base tree, built as its own
  \path{TreeClass::Big} instance --- the ``no augmentation'' reference point. \\
\path{KiB8} & 251 & 8\,KiB & Only $\sim$1.2--1.3$\times$ root-restart reduction ---
  barely more than off. \\
\path{KiB16} & 507 & 16\,KiB & $\sim$2$\times$/$\sim$1.6$\times$ reduction
  (district/warehouse). \\
\path{KiB32} (\textbf{default}) & 1019 & 32\,KiB & OLAP scan throughput fully
  recovered ($\sim$100\% of the untouched-leaf baseline) while root restarts are
  still down $\sim$4.4$\times$/$\sim$2.6$\times$ --- the measured sweet spot. \\
\path{KiB64} & 2043 & 64\,KiB & $\sim$8$\times$/$\sim$4.6$\times$ reduction, at a
  real but moderate OLAP scan cost ($\sim$82\% of baseline). \\
\path{KiB512} & 16379 & 512\,KiB & Maximum measured contention reduction
  ($\sim$32$\times$/$\sim$18$\times$), at a steep OLAP scan cost ($\sim$21\% of
  baseline). \\
\bottomrule
\end{tabular}
\end{center}

\textbf{The upper bound's mechanism is unavoidable, not a tuning oversight:} a
leaf's \path{overflow\_records\_count() == NUM\_RECORDS} ceiling is also
\emph{exactly when a leaf's dead/superseded versions get compacted away}, and
\path{RangeQueryIter} scans a leaf's entire physical record array, live and dead
alike. A bigger \path{NUM\_RECORDS} therefore means more accumulated dead-version
garbage sits in the leaf between compactions, and every range scan over that table
pays for wading through it. There is no way to decouple ``avoid frequent
compactions'' from ``keep dead-version garbage bounded'' in this design --- they are
the same operation, which is exactly why \path{BigTreeSize} exists as an explicit,
chosen point rather than a single hand-picked number.

\textbf{The lower bound is a proven, not merely observed, floor.} Every split-push
into an internal node \emph{appends} two fresh entries rather than overwriting
(\path{on\_overflow\_node}); the superseded entry goes dead until GC compacts it.
A freshly split node (two live children) needs \path{FAN\_OUT >= 4} just to have
room for one further push, and each subsequent push costs two more raw slots against
a budget of \path{FAN\_OUT - 1}. At \path{FAN\_OUT = 3} that margin is exactly zero:
a fresh root sits at the overflow threshold permanently, before it can accept even
one push --- a genuine, concurrency-\emph{independent} deadlock, confirmed by tracing
through \path{lacks\_room\_for\_split\_entries}/\path{on\_overflow\_node}/
\path{unsafe\_degree\_root}. Larger-but-still-small \path{FAN\_OUT} compiles and is
not provably deadlocked the same way, but empirically degrades from ``fine'' to
``indistinguishable from hung'' as concurrency rises: \path{FAN\_OUT} = 5, 7, 9, and
11 all ran cleanly at 1 warehouse/2 terminals but never completed (150s+ wall time
for a 20s run, under 15s of actual CPU work across all threads --- a genuine stall,
not slow-but-live progress) at 2 warehouses/4 terminals. \path{FAN\_OUT = 27}
(\path{KiB1}) was the smallest value confirmed to hold up at that load and at 4
warehouses/8 terminals; nothing here was tested past 8 terminals, so treat it as
``more contention-sensitive than \path{KiB2}/\path{KiB4}'', not as unconditionally
safe at arbitrary concurrency.

\subsection{A benchmark methodology bug, found and corrected mid-investigation}

An early benchmark (8 warehouses, 8s/point, 20 points across 5 sizes $\times$ 2
thread counts $\times$ gc on/off, on a 2$\times$ AMD EPYC 7742 server) found scan p50
latency essentially flat across all five then-existing size variants
(Tiny/Small/Medium/Large/Huge --- the naming predating today's \path{KiB\emph{N}}
scheme) and concluded ``Medium ties for best-or-tied-for-best on every metric
measured.'' That flatness finding was \textbf{wrong, as an artifact of the benchmark's
own duration}, not because the underlying effect does not exist. The engine's
always-on OLAP thread ages its scan snapshot by a configurable delay before
scanning; the original 8-second points never let that delay reach more than
3--4 seconds, so the garbage-accumulation effect the doc comment above describes
never actually got exercised within the timed window.

Re-run with a 30-second window and the \path{sweep} OLAP mode (isolating exactly the
warehouse+district scan), \path{olap\_param=5} (delays 0,1,2,3,4,5s):

\begin{center}
\begin{tabular}{lrr}
\toprule
Delay & \path{KiB32} (default) & \path{KiB512} \\
\midrule
0\,s & 155\,$\mu$s & 159\,$\mu$s \\
1\,s & 56\,$\mu$s & 249\,$\mu$s \\
2\,s & 39\,$\mu$s & 294\,$\mu$s \\
3\,s & 58\,$\mu$s & 343\,$\mu$s \\
4\,s & 61\,$\mu$s & 354\,$\mu$s \\
5\,s & 60\,$\mu$s & 269\,$\mu$s \\
\bottomrule
\end{tabular}
\end{center}

Both scans return the identical 88 rows every time --- the latency gap is pure
physical-record-touched overhead, not output size. \path{KiB512} is consistently
4--6$\times$ slower than \path{KiB32} past the first (near-zero-garbage) point,
confirming the original ``$\sim$100\%$\to$$\sim$21\% of baseline OLAP throughput''
claim in the source comment and correcting the earlier ``flat'' finding: it was an
artifact of too-short aging, not an absent effect. The lesson generalizes past this
one benchmark: \emph{a benchmark whose measured window is shorter than the
mechanism's own time constant will report ``no effect'' regardless of whether one
exists} --- exactly the caution the catalog documents already give about not
trusting a short/noisy sample, applied here to duration rather than run count.

A system-wide diagnostic added during the same investigation
(\path{mv\_test::SCAN\_TRACE}, off by default) found \textbf{5.74$\times$
records-visited per record-matched}, identical at both leaf sizes tested --- expected,
since the other twelve standard-sized tables' read volume swamps warehouse/district's
share of the global counter, but a real, general finding worth keeping in mind: this
engine's reads, at this write-heavy TPC-C scale, touch 5--6$\times$ more physical
records than they return, on average, across the whole workload, not just this one
optimization's target tables.

\subsection{Recommendation}

\path{KiB32} remains the correct general-purpose default: it is the current
\path{\#[default]} variant, and per the doc-commented measurements above it captures
most of the achievable contention win while keeping OLAP scan cost at parity with the
untouched baseline. Reaching for \path{KiB64}/\path{KiB512} is a deliberate trade of
scan latency for root-contention headroom that should be a conscious choice for a
specific deployment's read/write mix, not a default.

\section{Rejected, reverted, or deliberately optional experiments}
\label{sec:rejected}

An accurate account of this implementation has to include what did not survive.
Some of the entries below are already referenced above; this section collects all
of them, plus two GC-related attempts that are not documented in any other file in
this repository and exist only in this investigation's own history.

\begin{longtable}{IORGKS}
\toprule
\textbf{Location} & \textbf{Attempt} & \textbf{Why attempted} & \textbf{Observed result} &
\textbf{Failure mode} & \textbf{State} \\
\midrule
\endhead
\commit{d565b2d}--\commit{5c26abb} & ORWC optimistic read/write coupling & Non-blocking
  upgradeable writers alongside free readers. & $\sim$460 lines ultimately removed,
  no retained gain. & Repeated upgrade races and dangling readers made it duplicate
  OLC without a sound distinct guarantee. & \statereverted \\
\commit{45cb048} & Hold one long-lived range-scan block guard & Avoid re-borrowing
  blocks during iteration. & Not retained. & GC did not know about the guard's
  lifetime; reuse could invalidate it. & \statereverted \\
\S\ref{sec:scans} rejected set & Eager child-list reuse, small-page dedup, hybrid
  traversal, lean frames & Remove allocation/hash work, release internal pages early. &
  An experimental version measured 18--20\% faster. & Equal-boundary dedup cannot
  prove historical intervals superseded; GC stress returned duplicate keys. &
  \statereverted \\
\commit{0659d34} & Retry an alternate merge candidate immediately & Increase merge
  success on first-candidate failure. & No gain retained. & Tight retry loop lacked
  backoff, starved workers under contention. & \statereverted \\
Companion to \commit{85ac1b0} & Lower the leaf overflow trigger for more compaction
  headroom & Reduce chance of a full leaf meeting a pending write. & Stress hung. &
  More SMOs amplified contention while the actual off-by-one bug remained
  unaddressed. & \statereverted \\
\path{tx\_context.rs}/\path{smo.rs}, per-reader LCB check & Replace the global
  oldest-snapshot watermark with a per-reader liveness check
  (\path{any\_live\_reader\_missing}) & Avoid coupling every leaf's reclaim to the
  single slowest live transaction anywhere. & Instrumented run: 517{,}748 blocked
  deletes in 12 seconds, one fixed reader never advancing, no progress. & Since an
  open transaction may read any key, ``protect in case this reader needs it'' is
  indistinguishable in the worst case from protecting nearly everything --- a
  self-reinforcing compaction livelock, just computed more expensively than the
  global-watermark version it replaced. & \statereverted \\
\path{smo.rs}, \path{unsafe\_degree\_root} & Trip \code{Overflow} early once a
  leaf-root's dead-record count alone reaches 50\% of capacity (not only at 100\%
  physical fullness) & Reduce OLAP scan cost on hot tiny tables (see
  \S\ref{sec:bigtree}) without waiting for full-capacity compaction. & Scan latency
  win confirmed the mechanism (\path{KiB512}: $\sim$250--350\,$\mu$s
  $\to$ $\sim$145--195\,$\mu$s), but \code{cargo test --release} run repeatedly
  surfaced an intermittent ($\sim$30--40\% of runs) \path{stock s\_ytd} invariant
  failure, single-threaded, reproduced in isolation. & An early, ratio-based
  compaction trigger is reachable at a much lower bar than the existing 100\%-full
  trigger, in a regime (a same-transaction write to a leaf-root replaced mid-
  transaction by a different key's compaction) the surrounding abort/write-set
  machinery was never proven against. & \statereverted \\
Thread-local raw lock-free WAL buffer & Reuse encoding memory on the caller thread &
  Avoid one allocation per write. & No production gain. & Every real caller already
  routes through \path{LocalBatch}'s persistent buffer, even at batch size 1. &
  \statereverted \\
\commit{6486d97}, mimalloc feature & Cheaper cross-thread frees for the sharded WAL &
  YCSB-A +2.3--3.0\%; TPC-C fell 3.3--4.9\% across every backend. & Workload-dependent
  allocator trade-off. & Not a universal win. & \stateoptional \\
Very large \path{BigTreeSize} leaves & Defer compaction on the hottest tiny tables
  aggressively & Short tpmC results were flat and restart counts kept improving with
  size. & Aged scans at \path{KiB512} were 4--6$\times$ slower than \path{KiB32}. &
  Diminishing/inverted returns past the measured sweet spot. & \stateoptional \\
\bottomrule
\caption{Rejected, reverted, or deliberately optional experiments.}
\end{longtable}

\section{Open issues, as of this writing}
\label{sec:open}

These are real, unresolved gaps surfaced by this investigation's own history. They
are reported here, distinctly from the reverted-experiments table above, because
each one is a currently-\emph{unfixed defect} rather than a tried-and-abandoned idea
--- an honest account of the implementation has to include both categories
separately.

\begin{itemize}
  \item \textbf{A pre-existing \code{split()} livelock, independent of the two
        reverted GC attempts above.} While batch-verifying an unrelated fix,
        \path{concurrent_multi_step_transactions_across_a_small_shared_key_set_never_lose_an_update}
        hung for 500+ seconds (expected $\sim$3--4s) with steadily growing RSS at
        only $\sim$3\% CPU --- reproduced on \emph{completely unmodified} \texttt{HEAD}
        at roughly 1/25 runs ($\sim$4\%). Some data/version distribution can make
        \path{split()}'s \path{ByVersion} retry loop never converge, matching that
        function's own doc-commented predicted failure mode
        (``repeats forever''). A real fix needs either (a) a fallback to a genuine
        physical \path{ByKey} split when \path{ByVersion} compaction fails to shrink
        the survivor set enough, or (b) per-transaction read-set tracking so
        protection scopes to keys a reader actually touched rather than everything
        after its \path{ts\_start}. Both are materially larger than a single-predicate
        change. \stateopen{}
  \item \textbf{A second, independent GC-visibility gap in \code{record\_survives\_gc}.}
        OSIC visibility is LCB-based, not a raw \path{ts_start} comparison: a
        reader's own \path{ts_start} being numerically later than a foreign write's
        does not guarantee that write is visible to the reader yet ---
        \path{is_snapshot_live} can go false (the writer committed and unregistered)
        while some \emph{other}, still-registered reader's snapshot still needs to
        resolve to that exact superseded version. Both attempted fixes for this
        (the global-watermark version and the per-reader LCB version, described in
        \S\ref{sec:rejected}) were reverted because they reproduce a self-reinforcing
        compaction livelock. The gap itself remains --- documented in-code as a
        comment, with no behavior change --- and is the confirmed root cause of a
        residual $\sim$7--15\%-under-load flake in
        \path{verify_concurrent_shared_keys} (diagnosed with hard instrumentation
        evidence: the only surviving record for a missed key was a fresh insert not
        yet visible to the reader, with the previously-visible version already gone,
        immediately after a burst of \path{ByVersion} compactions on that leaf).
        \stateopen{}
  \item \textbf{TPC-C's OCC conflict rate is very high} at scale (Payment
        $\sim$71--76\% of attempts conflict and are discarded; New-Order
        $\sim$47\%, at 8 warehouses/64 terminals) --- pure wasted traversal/split
        work, unrelated to and unaddressed by the WAL work in \S\ref{sec:osic}.
        \stateopen{}
  \item \textbf{Contention split and XMerge} (Alhomssi \& Leis's B-tree
        contention/space-management techniques, from the same paper OSIC comes from)
        remain unported. Restart-attribution instrumentation run at this project's
        default TPC-C scale found root-level contention dominates (before the
        \code{BigTreeSize} fix); after that fix, the remaining restarts are mostly
        internal/routing-page (\code{checked\_live\_version}) failures on a small
        number of high-fan-in pages with hundreds of distinct keys funneling through
        each --- exactly the ``many distinct keys sharing one page's latch'' pattern
        contention split targets. No per-page restart/contention counters exist yet
        to drive it. \stateopen{}
  \item \textbf{The per-write WAL allocation on \code{WalWriter}'s cross-thread path}
        is smaller after the framing-hint fix (\S\ref{sec:osic}) but not eliminated:
        a fresh buffer is still allocated per write and freed on a different thread.
        A bounded per-shard buffer pool, with buffers returned right after
        \path{flush_loop} copies their bytes (well before the \path{pwrite}/
        \path{fsync}, so the pool need not wait on durability), is the identified but
        unimplemented fix. \stateopen{}
\end{itemize}

\section{Fresh measurements, this session}
\label{sec:fresh}

Everything in this section was run on the machine this document was written on
(AMD Ryzen 7 2700X, 8 cores/16 threads, rustc/cargo 1.97.1) against the exact
repository state at commit \commit{64e4380} (2026-08-12), specifically to
cross-check the historical claims above against the code as it exists today, not to
replace them --- the historical numbers above were mostly measured on different,
often larger/newer, hardware, and are kept as-is with their original attribution.

\subsection{Mechanism presence check}

Every mechanism this document describes as currently active was confirmed present
in the source at \commit{64e4380} by direct inspection/grep, not assumed from the
documents alone: the self-overwrite fast path (\path{insertion_stamp() == stamp} in
\path{insert\_on\_tree}/\path{update\_on\_tree}), \path{BigTreeSize}/\path{TreeClass}
(now an eight-variant enum, richer than any prior document describes ---
\S\ref{sec:bigtree} reflects the current, not the originally-documented, version),
\path{RESTART\_TRACE} (\path{false} by default, with its OOM-fix reset call intact),
\path{EXP\_SPIN\_CAP}/\path{JITTER\_BACKOFF\_THRESHOLD} in the SmartCell backoff
schedule, \path{LockFreeWalWriter}, and the point-read fix (no \path{skip\_while}
shortcut remains in \path{key\_point\_read\_from\_root}, now committed as
\commit{c0fe172} --- it was still uncommitted as of the most recent prior
investigation notes available). The two GC-livelock revert attempts documented in
\S\ref{sec:rejected} were confirmed genuinely absent from the current source
(\path{any\_live\_reader\_missing}/\path{is\_visible\_nocache} do not appear
anywhere in \path{src/}), consistent with those reverts having held.

\subsection{Full test suite}

\code{cargo test --release} at \commit{64e4380}: \textbf{151 passed, 0 failed, 0
ignored, finished in 5.71s} (wall clock for the whole invocation including the two
empty loom-test binaries: 6.14s). This is a healthy, fast-feedback current state,
consistent with the ``smoke-scale by default, \code{CMVBT\_FULL\_BENCH=1} for full
scale'' design described in \S\ref{sec:workload}. It does not, by design, exercise
either open livelock in \S\ref{sec:open} deterministically --- both are reported as
load-sensitive/rare, and this single run neither reproduces nor disproves them.

\subsection{Leaf layout benchmark, independently reproduced}

Following the exact reproduction recipe in \S\ref{sec:physical} (\code{rustc
--edition=2024 -C opt-level=3 -C target-cpu=native tools/leaf\_layout\_bench.rs},
\code{taskset -c 0}), five full runs of the binary were taken (the benchmark
internally already samples for 750\,ms per operation; five \emph{process} runs were
taken to match the original report's ``median of five'' protocol at the process
level). The five runs agreed with each other to well under 1\% --- this machine was
otherwise idle, so the tight agreement reflects a quiet machine, not that the
underlying effect is small.

\begin{center}
\begin{tabular}{lrrr}
\toprule
Operation (62.5\,MiB multi-leaf sweep) & AoS & Split & Speedup \\
\midrule
Point: latest visible & 10.90\,Mops/s & 11.27\,Mops/s & 1.033$\times$ \\
Point: 1-version fallback & 7.86\,Mops/s & 8.43\,Mops/s & 1.073$\times$ \\
Point: 3-version fallback & 4.66\,Mops/s & 4.32\,Mops/s & 0.927$\times$ \\
Scan: full leaves & 245.2\,Mslots/s & 236.2\,Mslots/s & 0.963$\times$ \\
Scan: 50\% ranges & 244.5\,Mslots/s & 237.3\,Mslots/s & 0.971$\times$ \\
Scan: 20\% ranges & 254.0\,Mslots/s & 236.0\,Mslots/s & 0.929$\times$ \\
\bottomrule
\end{tabular}
\end{center}

\textbf{This both confirms and diverges from the original report.} The direction of
the point-lookup effect replicates (split is faster for the common ``latest visible''
and ``one fallback'' cases), but at roughly a third of the claimed magnitude, and the
3-fallback case and two of three scan variants \emph{reverse sign} entirely versus
the original Ryzen 9 5900X report. The most plausible explanation is genuinely the
hardware, not a methodology error: this machine's CPU (Ryzen 7 2700X, Zen+, 2018) has
a smaller, higher-latency L3 cache and weaker prefetchers than the original report's
Ryzen 9 5900X (Zen 3, 2020) --- the split layout's win is specifically a
cache-pressure/locality effect (fewer unrelated cache lines touched per rejected
key), so a CPU with less cache to relieve pressure on, or a memory subsystem that
already dominates the access cost regardless of layout, would show a smaller and
noisier version of the same effect, plausibly crossing zero for the scan case where
the layout-independent cost (touching every byte) already dominates. This is
reported in full, with the divergence unresolved rather than explained away, because
a microbenchmark result that fails to reproduce cleanly on different hardware is
itself exactly the kind of finding \S\ref{sec:physical}'s own ``Reading the
performance claims'' caution warns a reader to expect and account for.

\subsection{WAL writer throughput, independently reproduced}

\code{cargo test --release} of \path{compare\_wal\_writer\_vs\_lockfree\_throughput}
(the same test the original synthetic comparison in \S\ref{sec:osic} used, but here
run in a full release+LTO build rather than debug, with a fuller thread sweep
$\{1,2,4,8,16\}$ and batch-size sweep $\{1,4,16,64,256\}$, 20{,}000 ops/thread):

\begin{center}
\begin{tabular}{lrrr}
\toprule
Backend & @1 thread & @8 threads & @16 threads \\
\midrule
\path{WalWriter} (4-way sharded channel) & 700{,}592 ops/s & 2{,}607{,}583 ops/s &
  4{,}339{,}121 ops/s \\
\path{LockFreeWalWriter}, unbatched & 384{,}583 ops/s & 839{,}536 ops/s &
  763{,}438 ops/s \\
\path{LockFreeWalWriter} + \path{LocalBatch(4)} & 1{,}021{,}650 ops/s &
  2{,}371{,}383 ops/s & 2{,}459{,}184 ops/s \\
\path{LockFreeWalWriter} + \path{LocalBatch(16)} & 1{,}509{,}518 ops/s &
  4{,}898{,}065 ops/s & 5{,}563{,}703 ops/s \\
\path{LockFreeWalWriter} + \path{LocalBatch(64)} & 1{,}475{,}161 ops/s &
  5{,}979{,}338 ops/s & 7{,}578{,}144 ops/s \\
\path{LockFreeWalWriter} + \path{LocalBatch(256)} & 1{,}515{,}828 ops/s &
  6{,}386{,}070 ops/s & \textbf{8{,}678{,}745 ops/s} \\
\bottomrule
\end{tabular}
\end{center}

This reproduces the qualitative story in \S\ref{sec:osic} cleanly: unbatched
lock-free writes lose to the sharded channel past a handful of threads, and
per-thread batching is what actually wins, monotonically improving with batch size
at 16 threads. The absolute numbers are considerably higher than the original
debug-build report (e.g.\ sharded-channel at 16 threads: 4.34\,M ops/s here vs.\ the
original 1.35\,M ops/s; batch-64 lock-free: 7.58\,M ops/s here vs.\ 5.09\,M ops/s
originally) --- expected and not a discrepancy worth worrying about, since this run
used a full \code{opt-level=3}/\code{lto=fat}/\code{codegen-units=1} release build
(this crate's actual \path{[profile.release]}) against the original's stated debug
build. The \emph{ranking} between backends is what the original report's conclusion
rested on, and that ranking reproduced exactly, magnitudes aside.

\section{Combined measured gains at a glance}
\label{sec:combined}

This table repeats only controlled, repository- or session-backed quantitative
results, extended from the earlier catalog's own closing table with this session's
fresh reproductions. It is the correct table to cite when a numerical gain is
required; do not attribute a combined-batch row to one component in isolation.

\begin{longtable}{p{5.6cm}p{4.4cm}p{4.6cm}p{4.0cm}}
\toprule
\textbf{Optimization set} & \textbf{Workload} & \textbf{Measured change} &
\textbf{Qualification} \\
\midrule
\endhead
4\,KiB base-block tuning (\commit{b29d347}) & Allocation / TLB profile & Alignment
  25\%$\to$100\%; dTLB load misses $-9$ to $-12$\% & Exact object layout dependent. \\
CRC32 + effective native build + 4-way WAL sharding & TPC-C, 16/64 threads &
  +7.4\% / +17.6\% & Combined result; do not assign all gain to one component. \\
Same three WAL/build changes & YCSB-A, 16/64 & +25.4\% / +68.7\% & Write-heavy
  Zipfian workload. \\
Same three WAL/build changes & YCSB-E, 16/64 & +16.2\% / +49.8\% & Scan-heavy with a
  write component. \\
Same three WAL/build changes & YCSB-F, 16/64 & +20.7\% / +70.1\% & Read-modify-write. \\
Ticket-less WAL + exact buffer hints & TPC-C / YCSB-A & +12.6\% / +25.3\% & Same
  short profiling setup, both fixes together. \\
Ticket-less WAL + exact buffer hints & Synthetic \path{WalWriter}, 16 threads &
  +79\% & 8-byte payload makes the removed ticket allocation unusually dominant. \\
Lock-free local batch-64 (original, debug build) & Synthetic WAL, 16 threads &
  5.09M vs.\ 1.35M ops/s & Unbatched lock-free is slower; see this session's release
  reproduction below for the same ranking at different absolute magnitudes. \\
\textbf{This session:} lock-free local batch-256 vs.\ sharded channel, release+LTO &
  Synthetic WAL, 16 threads & 8.68M vs.\ 4.34M ops/s (unbatched lock-free: 0.76M) &
  Same ranking as the original debug-build result; magnitudes not directly
  comparable across build profiles. \\
Zero-copy streaming range terminal & 200k-row full scan & 1.165--1.173$\times$;
  $\sim$+16.8\% & 40 alternating samples after warmup. \\
Range-before-visibility predicate & CH Q4 / Q5 & $\sim$+18\% / $\sim$+11--13\% &
  Indicative short contended runs; narrow ranges only. \\
Specialized count terminal & Full scan & $\sim$49--53M rows/s & Absolute
  throughput, retained-implementation number only. \\
\code{BigTreeSize KiB32} default vs.\ untouched leaf & District/warehouse root
  restarts & $\sim$4.4$\times$/$\sim$2.6$\times$ reduction & OLAP scan throughput at
  parity with baseline at this size. \\
\code{BigTreeSize KiB512} vs.\ \code{KiB32} & Aged OLAP scan latency (30s window,
  \code{sweep} mode) & 4--6$\times$ slower past 1s of aging & Corrects an earlier
  8s-window benchmark that read as ``flat.'' \\
Full-row YCSB-A batch & YCSB-A & 2.52M$\to$3.11M ops/s (+23.5\%) & 200k rows, 4
  threads, 3s. \\
Standard single-field + alias Zipf + hit-only path & YCSB-A (standard) / YCSB-C &
  4.81M ops/s / 7.16M$\to$14.82M ops/s (+106.8\%) & Different semantics than
  full-row mode; not directly comparable to it. \\
Scan-loop and measurement batch & YCSB-E & 4.59M$\to$6.22M ops/s (+35.4\%) &
  Combined batch result; latency samples one scan in 1,024. \\
Mimalloc feature & YCSB-A / TPC-C & +2.3--3.0\% / $-3.3$ to $-4.9$\% & Kept optional,
  not a universal optimization. \\
\textbf{This session:} leaf layout, split vs.\ AoS, Ryzen 7 2700X & Point lookup /
  scan, 62.5\,MiB multi-leaf & 1.03--1.07$\times$ / 0.93--0.97$\times$ & Smaller and
  partly reversed vs.\ the original Ryzen 9 5900X report --- see \S\ref{sec:fresh}
  for the hardware-sensitivity discussion. \\
\textbf{This session:} full test suite, commit \commit{64e4380} & Development test
  run & 151 passed, 0 failed, 5.71s & Correctness/operability result, current as of
  this document's writing. \\
\bottomrule
\end{longtable}

\section{Timeline}
\label{sec:timeline}

Table~\ref{tab:timeline} indexes the commits discussed in this document
chronologically, for cross-reference against \code{git log}. It is provided as a
navigational aid only --- the organizing structure of this document is architectural
layer and dependency, per \S\ref{sec:intro}, not this chronology.

\begin{longtable}{@{}p{2.3cm}p{1.9cm}p{9.2cm}@{}}
\toprule
\textbf{Commit} & \textbf{Date} & \textbf{Theme} \\
\midrule
\endhead
\commit{acdc7a8} & 2024-02-14 & GC started (\S\ref{sec:gc}) \\
\commit{95e149f} & 2024-02-29 & GC active-tx index moved to a concurrent B+-tree (\S\ref{sec:gc}) \\
\commit{d565b2d}--\commit{5c26abb} & 2024-01--2024-12 & ORWC implemented, iterated, deprecated, purged (\S\ref{sec:cc}) \\
\commit{48c4ded}, \commit{a7bb29c}, \commit{a518eed} & 2024-05--2025-04 & Fence removal / ordering tightening (\S\ref{sec:cc}) \\
\commit{bea1117} & 2025-06-05 & GC-driven update-in-place (unrelated mechanism, \S\ref{sec:selfoverwrite}) \\
\commit{17b8f5f}--\commit{836a1db} & 2025-12--2026-01 & Pre-OSIC commit/visibility, thread-watermark patching (\S\ref{sec:osic}) \\
\commit{34555e2} & 2026-06-20 & Overflow-trigger tuning (\S\ref{sec:smo}) \\
\commit{2a2be93}, \commit{40b00fd} & 2026-06 & Further ordering relaxation (\S\ref{sec:cc}) \\
\commit{714a82f} & 2026-07-10 & OSIC/WAL adopted wholesale (\S\ref{sec:osic}) \\
\commit{3e8b837}, \commit{ab10660}, \commit{73cc4a5} & 2026-07-13/15 & Group commit, ELR removal, log truncation (\S\ref{sec:osic}) \\
\commit{5f92e9f} & 2026-07-15 & Tables-as-shards (\S\ref{sec:sharding}) \\
\commit{7f2bd09} & 2026-05-21 & Cache-line isolation for node tag (\S\ref{sec:cc}) \\
\commit{45c7c21}/\commit{fd92f38}, \commit{d0e0977} & 2026-07-21 & Conservative child latch removed; retired-bit race fix (\S\ref{sec:cc}, \S\ref{sec:gc}) \\
\commit{f9fc960} & 2026-07-23 & GC own-shard-first + steal (\S\ref{sec:sharding}) \\
\commit{ecc9900} & 2026-07-27 & Racey-stamp + dead-page-key collision fixes (\S\ref{sec:gc}) \\
\commit{fe03d90} & 2026-07-28 & Obsolete-mark removed; overlap-based liveness (\S\ref{sec:gc}) \\
\commit{4bc78de} & 2026-07-28 & Merge-overflow accounting fix (\S\ref{sec:smo}) \\
\commit{7432525}--\commit{fc8806b} & 2026-07-31 & TPC-C crash hardening week (\S\ref{sec:smo}) \\
\commit{160186e}, \commit{85ac1b0}, \commit{fbedd74} & 2026-08-01 & Repeated-key abort fix; \textbf{self-overwrite fast path}; hot-key repro (\S\ref{sec:selfoverwrite}) \\
\commit{2a2607e}, \commit{dc30111}, \commit{dcb5abe}, \commit{0659d34} & 2026-08-02 & Split off-by-one and merge/boundary-search fixes (\S\ref{sec:smo}) \\
\commit{eac6e40}, \commit{b29d347} & 2026-08-02 & Cache-padded worker atomics; FAN\_OUT page alignment (\S\ref{sec:sharding}, \S\ref{sec:smo}) \\
\commit{4a1a50a} & 2026-08-09 & WAL channel sharding, table-driven CRC32 (\S\ref{sec:osic}) \\
\commit{ebdd789}, \commit{9ced1a3}, \commit{efc3773} & 2026-08-10 & Lock-free WAL backend; RESTART\_TRACE OOM; cross-shard race; ticket/framing fixes (\S\ref{sec:osic}) \\
\commit{8e48ae5} & 2026-08-11 & Range-before-visibility reorder; \code{dyn} removal (\S\ref{sec:scans}) \\
\commit{c0fe172} & 2026-08-11/12 & Point/range-read miss fix: physical order does not track \path{ts\_start} order (\S\ref{sec:scans}) \\
\commit{1f6d666}, \commit{c525eba} & 2026-08-09/12 & \code{BigTreeSize} introduced, then extended with sub-4KiB variants (\S\ref{sec:bigtree}) \\
\bottomrule
\caption{Chronological index of commits discussed in this document.}
\label{tab:timeline}
\end{longtable}

\section{Conclusion}
\label{sec:conclusion}

Read end to end, the history of this codebase tells one consistent story at every
layer: an initial design gets something working with coarse-grained locks, mutexes,
or bolted-on atomics/fences; profiling or ThreadSanitizer then reveals exactly where
that coarseness costs correctness (a race) or throughput (contention); and the fix
is, more often than not, not ``add a lock'' but ``restructure the data so the hazard
cannot occur'' --- append-only pages instead of in-place edits, interval-overlap-derived
liveness instead of a mutable obsolete bit, per-worker sharded bookkeeping instead of
one shared structure, instant in-memory commit instead of commit-time fsync. The
four foundational optimizations named in \S\ref{sec:foundational} are the load-bearing
instances of that pattern; the self-overwrite trick in \S\ref{sec:selfoverwrite} is a
small, local instance of exactly the same shape one layer up: rather than adding
machinery to garbage-collect or specially handle a hot key's pile of uncommitted
versions, the fix removes the pile-up from occurring in the first place.

An equally important, less flattering thread runs alongside it, and this document
has tried not to understate it: not every attempted instance of this pattern
succeeded. Two separate attempts to close the GC-visibility gap in
\S\ref{sec:rejected} each independently rediscovered the same self-reinforcing
compaction livelock from a different angle, and both were correctly reverted rather
than shipped with a rare, hard-to-attribute correctness cost. A benchmark's own
measured window can itself be wrong in a way that produces a confident, false
``no effect'' conclusion (\S\ref{sec:bigtree}), and a microbenchmark's headline
percentage can fail to reproduce, in both magnitude and sign, on different hardware
without either measurement being in error (\S\ref{sec:fresh}). A codebase's real
state, at any given moment, is the sum of the optimizations that shipped, the ones
that were tried and correctly reverted, and the ones that are known to still be
broken --- this document has tried to give an accurate account of all three, not
just the first.

\begin{thebibliography}{9}

\bibitem{becker1996asymptotically}
Becker, B., Gschwind, S., Ohler, T., Seeger, B., Widmayer, P.
\emph{An asymptotically optimal multiversion B-tree}.
The VLDB Journal, 5(4):264--275, 1996.

\bibitem{tonta2026multiversion}
Tonta, A., Seeger, B., Soisalon-Soininen, E.
\emph{Multiversion Concurrency Control for Multiversion B-Trees}.
arXiv preprint arXiv:2606.09133, 2026.

\bibitem{alhomssi2023scalable}
Alhomssi, A., Leis, V.
\emph{Scalable and Robust Snapshot Isolation for High-Performance Storage Engines}.
Proceedings of the VLDB Endowment, 16(6):1426--1438, 2023.
doi:10.14778/3583140.3583157.

\end{thebibliography}

\end{document}
